
% !TeX program = xelatex
% 数学分析（二）期末复习讲义：扩写详解版
% 说明：本版封面仿 ElegantBook 风格：上半区图片 + 下半区文字信息。将 cover-top.jpg / cover-deco.png 放在同目录可自动替换。
\documentclass[UTF8,zihao=-4,openany]{ctexbook}

% -------------------- 页面与基础宏包 --------------------
\usepackage[a4paper,top=2.5cm,bottom=2.5cm,left=2cm,right=2cm]{geometry}
\usepackage{amsmath,amssymb,mathtools,amsthm,bm}
\usepackage{graphicx}
\usepackage{enumitem,booktabs,longtable,array,multicol}
\usepackage{xcolor}
\usepackage[most]{tcolorbox}
\usepackage{hyperref}
\usepackage{fancyhdr}
\usepackage{titlesec}
\usepackage{lastpage}
\usepackage{tikz}
\allowdisplaybreaks

% -------------------- 颜色与超链接 --------------------
\definecolor{mainblue}{HTML}{1F4E79}
\definecolor{lightblue}{HTML}{EAF3FB}
\definecolor{theogreen}{HTML}{2E7D32}
\definecolor{lightgreen}{HTML}{ECF7EC}
\definecolor{warnred}{HTML}{B71C1C}
\definecolor{lightred}{HTML}{FFF1F1}
\definecolor{softgray}{HTML}{F7F7F7}
\definecolor{deeporange}{HTML}{E65100}
\definecolor{covergray}{HTML}{F4F4F4}
\hypersetup{colorlinks=true,linkcolor=mainblue,urlcolor=mainblue,citecolor=mainblue}

% -------------------- 页面风格 --------------------
\linespread{1.20}
\setlength{\parindent}{2em}
\setlist[enumerate]{itemsep=0.35em,topsep=0.35em,leftmargin=2.2em}
\setlist[itemize]{itemsep=0.25em,topsep=0.25em,leftmargin=2em}
\setlength{\headheight}{14pt}
\pagestyle{fancy}
\fancyhf{}
\fancyhead[LE,RO]{\small\leftmark}
\fancyfoot[C]{\small\thepage\,/\,\pageref{LastPage}}
\renewcommand{\headrulewidth}{0.4pt}
\titleformat{\chapter}[hang]{\Huge\bfseries\color{mainblue}}{第\thechapter 章}{0.8em}{}
\titleformat{\section}[hang]{\Large\bfseries\color{mainblue}}{\thesection}{0.6em}{}
\titleformat{\subsection}[hang]{\large\bfseries\color{theogreen}}{\thesubsection}{0.6em}{}

% -------------------- 常用命令 --------------------
\newcommand{\R}{\mathbb{R}}
\newcommand{\N}{\mathbb{N}}
\newcommand{\dd}{\,\mathrm{d}}
\newcommand{\ee}{\mathrm{e}}
\newcommand{\ii}{\mathrm{i}}
\newcommand{\grad}{\nabla}
\newcommand{\diver}{\operatorname{div}}
\newcommand{\curl}{\operatorname{curl}}
\newcommand{\lap}{\Delta}
\newcommand{\rank}{\operatorname{rank}}
\newcommand{\sgn}{\operatorname{sgn}}
\newcommand{\norm}[1]{\left\lVert #1\right\rVert}
\newcommand{\abs}[1]{\left\lvert #1\right\rvert}
\newcommand{\coursekey}{\textcolor{deeporange}{\bfseries【课件重点】}}
\newcommand{\examkey}{\textcolor{warnred}{\bfseries【高频考点】}}
\newcommand{\thinking}{\textcolor{mainblue}{\bfseries【解题思路】}}
\newcommand{\point}{\textcolor{warnred}{\bfseries【考点】}}
\newcommand{\checklist}{\textcolor{theogreen}{\bfseries【检查清单】}}
\newcommand{\coverbanner}{cover-top.jpg}
\newcommand{\coverdeco}{cover-deco.png}
\newcommand{\coverauthor}{作者：Symphony}
\newcommand{\covertypesetter}{排版：GPT 5.5 thinking}
\newcommand{\coverschool}{课程：数学分析（二）}
\newcommand{\coverdate}{时间：\today}
\newcommand{\coverversion}{版本：Expanded 2.2}
\newcommand{\coverquote}{Analysis is the art of estimating the invisible.}

% -------------------- 盒子环境 --------------------
\tcbset{enhanced,breakable,boxrule=0.6pt,arc=2mm,left=1.5mm,right=1.5mm,top=1mm,bottom=1mm}
\newtcbtheorem[number within=chapter]{definition}{定义}{colback=lightblue,colframe=mainblue,fonttitle=\bfseries}{def}
\newtcbtheorem[number within=chapter]{theorem}{定理}{colback=lightgreen,colframe=theogreen,fonttitle=\bfseries}{thm}
\newtcbtheorem[number within=chapter]{lemma}{引理}{colback=lightgreen,colframe=theogreen,fonttitle=\bfseries}{lem}
\newtcbtheorem[number within=chapter]{corollary}{推论}{colback=lightgreen,colframe=theogreen,fonttitle=\bfseries}{cor}
\newtcbtheorem[number within=chapter]{example}{例题}{colback=softgray,colframe=black!55,fonttitle=\bfseries}{ex}
\newtcbtheorem[number within=chapter]{mistake}{易错点}{colback=lightred,colframe=warnred,fonttitle=\bfseries}{mis}
\newtcolorbox{methodbox}[1]{title={#1},colback=lightblue,colframe=mainblue,fonttitle=\bfseries}
\newtcolorbox{formula}[1]{title={#1},colback=white,colframe=theogreen,fonttitle=\bfseries}
\newtcolorbox{problem}[1]{title={#1},colback=white,colframe=mainblue,fonttitle=\bfseries}
\newtcolorbox{solutionbox}[1]{title={#1},colback=softgray,colframe=theogreen,fonttitle=\bfseries}
\newtcolorbox{warningbox}[1]{title={#1},colback=lightred,colframe=warnred,fonttitle=\bfseries}

\begin{document}
\frontmatter

% -------------------- ElegantBook 风格封面：上图下文，可替换图片 --------------------
\begin{titlepage}
\thispagestyle{empty}
\begin{tikzpicture}[remember picture,overlay]
  % 背景
  \fill[white] (current page.south west) rectangle (current page.north east);

  % 上半区大图：把你的主视觉图命名为 cover-top.jpg 放在 tex 同目录即可自动替换
  \IfFileExists{\coverbanner}{%
    \node[anchor=north,inner sep=0pt] at (current page.north){%
      \includegraphics[width=\paperwidth,height=0.60\paperheight,keepaspectratio]{\coverbanner}%
    };
  }{%
    % 没有图片时的极简占位：深色图片区，效果接近 ElegantBook 封面上方主视觉
    \fill[black!88] (current page.north west) rectangle ([yshift=-0.60\paperheight]current page.north east);
    \foreach \i in {1,...,90}{%
      \pgfmathsetmacro{\sx}{random()*\paperwidth}
      \pgfmathsetmacro{\sy}{random()*0.56*\paperheight}
      \fill[white,opacity=0.30] ([xshift=\sx pt,yshift=-\sy pt]current page.north west) circle (0.35pt);
    }
    \node[text=white!70,font=\Large,align=center] at ([yshift=-0.30\paperheight]current page.north) {cover-top.jpg\\主视觉图片占位};
  }

  % 下半区白色信息面板
  \fill[white] ([yshift=-0.60\paperheight]current page.north west) rectangle (current page.south east);

  % 标题与信息
  \node[anchor=west,align=left] at ([xshift=2.55cm,yshift=6.75cm]current page.south west) {%
    \begin{minipage}{0.64\paperwidth}
      {\zihao{1}\bfseries 数学分析（二）期末复习讲义}\par
      \vspace{0.65cm}
      {\zihao{4}\color{black!68} }\par
      \vspace{1.0cm}
      {\zihao{-4}\color{black!55}\coverauthor}\par\vspace{0.20cm}
      {\zihao{-4}\color{black!55}\covertypesetter}\par\vspace{0.20cm}
      {\zihao{-4}\color{black!55}\coverschool}\par\vspace{0.20cm}
      {\zihao{-4}\color{black!55}\coverdate}\par\vspace{0.20cm}
      {\zihao{-4}\color{black!55}\coverversion}\par
    \end{minipage}%
  };

  % 右下角装饰图：可选。把小插画命名为 cover-deco.png 即可自动出现
  \IfFileExists{\coverdeco}{%
    \node[anchor=south east,inner sep=0pt] at ([xshift=-2.35cm,yshift=2.45cm]current page.south east){%
      \includegraphics[width=4.1cm]{\coverdeco}%
    };
  }{%
    \draw[black!10,line width=0.8pt] ([xshift=-6.45cm,yshift=2.55cm]current page.south east) rectangle ++(4.15cm,3.2cm);
    \node[text=black!28,font=\small,align=center] at ([xshift=-4.38cm,yshift=4.15cm]current page.south east){cover-deco.png\\装饰图占位};
  }

  % 底部格言
  \node[anchor=south,align=center,text=black!85,font=\small\bfseries] at ([yshift=1.35cm]current page.south) {\coverquote};
\end{tikzpicture}
\vspace*{\fill}
\end{titlepage}

\chapter*{前言}
\addcontentsline{toc}{chapter}{前言}
\qquad 本讲义基于本学期课件压缩包、浙江大学数学分析（甲）II（H）2024--2025春夏回忆卷、中国科学技术大学数学分析A2期末资料、南开数学分析II历年试题汇总整理而成。正文按期末复习最需要的知识结构重写为六章：数项级数、函数项级数与一致收敛、重积分、曲线积分与曲面积分、含参变量积分、Fourier级数。正文部分有大部分定理的简略证明。结构为：核心定义、关键定理与证明脉络、典型解题方法、高频易错点(由于篇幅限制没有安排课后习题环节，可能在未来的版本中出现)\par
本讲义适合考前一周左右复习，若时间较为紧张的话也可以直接看例题的做题方法然后直接去做本校的往年试卷，本讲义选取的试卷事实上对基础的要求比较高，希望大家通过这本讲义都能学到更多。\par
鉴于本人水平有限，难免存在笔误和疏漏。如果大家在使用过程中发现错误，或者有更好的解题思路，欢迎加入数学协会交流群反馈：群号1033564584。本讲义的最新版与 \LaTeX{} 源码将持续更新于我的 GitHub 主页：
\href{https://github.com/Symphony666}{https://github.com/Symphony666}。
群内也会同步更新讲义的最新版本，和分享其他数学课程的复习资料和竞赛信息。


\tableofcontents
\mainmatter
\part{正文}

\chapter{数项级数}
\section{核心定义}
\begin{definition}{数项级数、部分和、收敛与发散}{series}
给定数列 $\{u_n\}$，形式和 $\sum_{n=1}^{\infty}u_n$ 称为数项级数。记
\[
S_n=u_1+u_2+\cdots+u_n.
\]
若部分和数列 $\{S_n\}$ 收敛于有限数 $S$，则称级数收敛，且称 $S$ 为级数的和；若 $\{S_n\}$ 不收敛，则称级数发散。注意：级数是否收敛，不是看“项越来越小”这一点，而是看“部分和是否趋于有限极限”。
\end{definition}

\begin{definition}{正项级数、绝对收敛与条件收敛}{absconv}
若 $u_n\ge 0$，则 $\sum u_n$ 称为正项级数。若 $\sum |u_n|$ 收敛，则 $\sum u_n$ 绝对收敛；若 $\sum u_n$ 收敛而 $\sum |u_n|$ 发散，则称为条件收敛。绝对收敛比普通收敛强，因为由Cauchy准则可知 $\sum |u_n|$ 收敛会推出任意尾和 $\sum_{k=m}^{n}u_k$ 的绝对值可被 $\sum_{k=m}^{n}|u_k|$ 控制。
\end{definition}

\section{关键定理与证明脉络}
\begin{theorem}{收敛必要条件与柯西准则}{cauchyseries}
若 $\sum u_n$ 收敛，则 $u_n\to0$。级数 $\sum u_n$ 收敛的充要条件为：任给 $\varepsilon>0$，存在 $N$，当 $n>m\ge N$ 时，
\[
\left|\sum_{k=m+1}^{n}u_k\right|<\varepsilon.
\]
\end{theorem}
\noindent\textbf{证明思路:} 级数收敛等价于部分和\(S_n\)收敛，而数列收敛的Cauchy准则正是 $|S_n-S_m|<\varepsilon$。由于 $S_n-S_m=\sum_{k=m+1}^{n}u_k$，便得到级数的尾和形式。必要条件由 $u_n=S_n-S_{n-1}\to S-S=0$ 得到。

\begin{theorem}{正项级数的核心判别法}{positiveseries}
设 $u_n,v_n\ge0$。
\begin{enumerate}
\item \textbf{单调有界原则：} 正项级数收敛当且仅当部分和数列有界。
\item \textbf{比较判别：} 若 $0\le u_n\le v_n$ 且 $\sum v_n$ 收敛，则 $\sum u_n$ 收敛；若 $0\le v_n\le u_n$ 且 $\sum v_n$ 发散，则 $\sum u_n$ 发散。
\item \textbf{极限比较：} 若 $\lim_{n\to\infty}u_n/v_n=c$ 且 $0<c<\infty$，则二者同敛散。
\item \textbf{比值判别：} 若 $\limsup |u_{n+1}/u_n|<1$，则绝对收敛；若 $\liminf |u_{n+1}/u_n|>1$，则发散。
\item \textbf{根值判别：} 若 $\limsup \sqrt[n]{|u_n|}<1$，则绝对收敛；若 $\limsup \sqrt[n]{|u_n|}>1$，则发散。
\item \textbf{积分判别：} 若 $f$ 在 $[1,\infty)$ 上非负、连续、单调下降，且 $u_n=f(n)$，则 $\sum u_n$ 与 $\int_1^\infty f(x)\dd x$ 同敛散。
\end{enumerate}
\end{theorem}

\noindent\textbf{证明思路:} \begin{enumerate}
    \item [1,]证明极限、比值、根值判别法只要简单地由极限的定义展开就可以证明。
    \item [2,]对于积分判别法，由条件(看图也很直观)我们有\[
    f(n+1)\leq \int_{n}^{n+1}f(x)\mathrm{d}x\leq f(n)
    \]累加求和就有\[
    \sum u_n\leq \int_{1}^{+\infty}f(x)\dd x\leq \sum u_n
    \]这就表明$\sum u_n$ 与 $\int_1^\infty f(x)\dd x$ 同敛散。
\end{enumerate}
\begin{theorem}{Leibniz判别法}{dirichletabel}
若 $a_n\downarrow0$\((\text{这里表示单调趋于0})\)，则 $\sum (-1)^{n-1}a_n$ 收敛。
\end{theorem}
\noindent\textbf{证明思路:}\\
令\(S_n=\sum_{k=1}^{n}(-1)^{n-1}a_n\)分奇偶之后用单调有界原理证明收敛，再证明收敛到同一个极限即可。

\begin{theorem}{abel判别法，dirichlet判别法}{abel-dirichlet-test}
    \begin{enumerate}
        \item[1] abel判别法：\(a_n\)单调有界，\(B_n\)收敛 \(\implies\sum a_nb_n\)收敛
        \item[2] dirichlet判别法：\(a_n\downarrow0\)，\(B_n\)有界 \(\implies\sum a_nb_n\)收敛
    \end{enumerate}
\end{theorem}
\textbf{证明思路：}Dirichlet与Abel的共同核心是Abel公式(这就是分部积分的离散形式)：
\[
\sum_{k=1}^{p}a_kb_k=a_pB_p-\sum_{k=1}^{p-1}(a_{k+1}-a_k)B_k
\]
只要 $B_k$ 有界且 $a_k$ 单调趋零，右端尾部就可被 $|B_m|$ 与 $\sum |a_{k+1}-a_{k}|$ 控制。考试中出现 $\sin nx,\cos nx,(-1)^n$ 等振荡因子时，往往要想到“部分和有界”。


\section{典型解题方法}
\begin{methodbox}{敛散性判断四步法}
\begin{enumerate}
\item \textbf{必要条件先行。} 先算 $u_n$ 是否趋于 $0$。若不趋于 $0$，级数必发散。
\item \textbf{识别主阶。} 对含根式、对数、三角函数的通项，先做等价无穷小或Taylor展开，找到与 $1/n^p$、$1/(n(\ln n)^q)$ 的关系。
\item \textbf{根据结构选工具。} 正项用比较、比值、根值、积分；阶乘/指数用比值；$n$ 次方根用根值；交错或振荡用Dirichlet/Abel。
\item \textbf{结尾补充绝对/条件。} 对一般项级数，先判 $\sum |u_n|$，再判原级数；不要只写“收敛”而漏掉条件收敛。
\end{enumerate}
\end{methodbox}

\begin{example}{}{}
判断 $\sum_{n=2}^{\infty}\dfrac{1}{n(\ln n)^p}$ 的敛散性。
\tcblower
由积分判别，考虑
\[
\int_2^\infty \frac{\dd x}{x(\ln x)^p}.
\]
令 $u=\ln x$，则 $\dd u=\dd x/x$，积分变为
\[
\int_{\ln2}^\infty \frac{\dd u}{u^p}.
\]
当 $p>1$ 时收敛，当 $p\le1$ 时发散。因此原级数当且仅当 $p>1$ 收敛。这个结论是期末常用比较对象。
\end{example}

\begin{example}{基本敛散性问题}{sum}
    \textbf{题目：} 判断级数 $\sum_{n=1}^\infty \left(n^{\frac{1}{n}} - 1\right)^p$ 的敛散性。
\tcblower
考虑先判断主阶：通项为幂指函数减1的形式，利用等价无穷小 $e^x - 1 \sim x \ (x \to 0)$ 降阶，转化为对数型级数后与p级数比较敛散性。\\
当 $n \to \infty$ 时，$n^{\frac{1}{n}} = e^{\frac{\ln n}{n}}$，且 $\frac{\ln n}{n} \to 0$。由等价无穷小替换：
$$
n^{\frac{1}{n}} - 1 = e^{\frac{\ln n}{n}} - 1 \sim \frac{\ln n}{n}
$$
因此通项的主阶为：
$$
\left(n^{\frac{1}{n}} - 1\right)^p \sim \frac{(\ln n)^p}{n^p}
$$

\begin{itemize}
    \item 当 $p > 1$ 时：取 $1 < r < p$，由 $\lim_{n\to\infty} \frac{(\ln n)^p}{n^{p-r}} = 0$，可知存在常数 $C>0$，使得 $n$ 充分大时 $\frac{(\ln n)^p}{n^p} \leq \frac{C}{n^r}$。而 $\sum \frac{1}{n^r}$ 收敛，由比较判别法，原级数收敛。
    \item 当 $p \leq 1$ 时：$n$ 充分大时 $\ln n > 1$，故 $\frac{(\ln n)^p}{n^p} \geq \frac{1}{n^p}$。而 $\sum \frac{1}{n^p}$ 发散，由比较判别法，原级数发散。
\end{itemize}
\end{example}

\begin{example}{}{}
判断级数 $\sum_{n=1}^\infty \frac{n!}{n^n}$ 的敛散性。
\tcblower
通项含阶乘 $n!$，是比值判别法的标志性题型。比值法适合处理相邻两项有明显约分关系的通项（阶乘、指数、幂次）。\\
设 $u_n = \frac{n!}{n^n} > 0$，计算相邻两项的比值：
$$
\frac{u_{n+1}}{u_n} = \frac{(n+1)!}{(n+1)^{n+1}} \cdot \frac{n^n}{n!} = \frac{(n+1)\cdot n^n}{(n+1)\cdot (n+1)^n} = \left( \frac{n}{n+1} \right)^n = \frac{1}{\left(1+\frac{1}{n}\right)^n}
$$

取极限：
$$
\lim_{n\to\infty} \frac{u_{n+1}}{u_n} = \lim_{n\to\infty} \frac{1}{\left(1+\frac{1}{n}\right)^n} = \frac{1}{e} < 1
$$

由正项级数的比值判别法，原级数收敛。
\end{example}

\begin{example}{}{}
讨论级数 $\sum_{n=2}^\infty (-1)^n \frac{(\ln n)^p}{n}$ 的敛散性（绝对收敛/条件收敛/发散）。
\tcblower
先验证必要条件，再判断绝对收敛性（正项级数），最后用莱布尼茨判别法判断条件收敛性。核心是对参数 $p$ 分类讨论。\\
设 $u_n = \frac{(\ln n)^p}{n}$，先分析正项级数 $\sum u_n$ 的敛散性（绝对收敛性）：
由积分判别法，考虑反常积分 $\int_2^{+\infty} \frac{(\ln x)^p}{x} dx$。令 $t = \ln x$，则 $dt = \frac{dx}{x}$，积分化为：
$$
\int_{\ln 2}^{+\infty} t^p dt
$$

\begin{itemize}
    \item 当 $p < -1$ 时：上述积分收敛，故正项级数 $\sum u_n$ 收敛，原级数\textbf{绝对收敛}。
    \item 当 $p \geq -1$ 时：上述积分发散，故正项级数 $\sum u_n$ 发散，原级数不绝对收敛。
\end{itemize}

接下来分析交错级数的条件收敛性（$p \geq -1$ 的情况）：
对交错级数 $\sum (-1)^n u_n$，验证莱布尼茨判别法的两个条件：
\begin{enumerate}
    \item 极限为0：$\lim_{n\to\infty} u_n = \lim_{n\to\infty} \frac{(\ln n)^p}{n} = 0$（对任意实数 $p$，对数增长速度慢于多项式）；
    \item 单调递减：令 $f(x) = \frac{(\ln x)^p}{x}$，求导得：
    $$
    f'(x) = \frac{p (\ln x)^{p-1} - (\ln x)^p}{x^2} = \frac{(\ln x)^{p-1}(p - \ln x)}{x^2}
    $$
    当 $x$ 充分大时，$\ln x > p$，故 $f'(x) < 0$，即数列 $\{u_n\}$ 单调递减。
\end{enumerate}

满足莱布尼茨条件，因此交错级数收敛。

综上结论：
\begin{itemize}
    \item $p < -1$：级数绝对收敛；
    \item $p \geq -1$：级数条件收敛。
\end{itemize}
\end{example}

\begin{example}{}{}
判断级数 $\sum_{n=1}^\infty \frac{1}{\int_0^n \sqrt{1+x^4} dx}$ 的敛散性。
\tcblower
先识别主阶：通项分母为积分形式，无需精确计算积分值，只需估计积分的无穷大阶数，找到分母与 $n^k$ 的等价关系，再用比较判别法。\\
先估计分母积分的阶：当 $x \to +\infty$ 时，$\sqrt{1+x^4} \sim x^2$，因此：
$$
\int_0^n \sqrt{1+x^4} dx \sim \int_0^n x^2 dx = \frac{n^3}{3} \quad (n \to \infty)
$$

因此通项的主阶为：
$$
\frac{1}{\int_0^n \sqrt{1+x^4} dx} \sim \frac{3}{n^3}
$$

而p级数 $\sum_{n=1}^\infty \frac{1}{n^3}$ 收敛（$p=3>1$），由比较判别法的极限形式，原级数收敛。
\end{example}

\begin{example}{}{}
讨论级数 $\sum_{n=3}^\infty \frac{1}{n \cdot (\ln n)^p \cdot (\ln \ln n)^q}$ 的敛散性。
\tcblower
积分判别法的经典进阶题型，通项正且单调递减，通过两次换元将积分逐步转化为p级数形式，是对数判别法的核心结论。\\
设 $f(x) = \frac{1}{x \cdot (\ln x)^p \cdot (\ln \ln x)^q}$，当 $x \geq 3$ 时，$f(x)$ 非负且单调递减，满足积分判别法条件。

考虑反常积分 $\int_3^{+\infty} f(x) dx$，做第一次换元：令 $u = \ln x$，则 $du = \frac{dx}{x}$，积分化为：
$$
\int_{\ln 3}^{+\infty} \frac{du}{u^p \cdot (\ln u)^q}
$$

分情况讨论：
\begin{enumerate}
    \item 当 $p > 1$ 时：
    取 $1 < r < p$，当 $u$ 充分大时，$\frac{1}{u^p (\ln u)^q} \leq \frac{1}{u^r}$。而 $\int^{+\infty} \frac{du}{u^r}$ 收敛，故原积分收敛，级数\textbf{收敛}。

    \item 当 $p = 1$ 时：
    做第二次换元：令 $v = \ln u$，则 $dv = \frac{du}{u}$，积分化为 $\int_{\ln \ln 3}^{+\infty} \frac{dv}{v^q}$。
    \begin{itemize}
        \item 若 $q > 1$：积分收敛，级数收敛；
        \item 若 $q \leq 1$：积分发散，级数发散。
    \end{itemize}

    \item 当 $p < 1$ 时：
    $u$ 充分大时，$\frac{1}{u^p (\ln u)^q} \geq \frac{1}{u}$，而 $\int^{+\infty} \frac{du}{u}$ 发散，故原积分发散，级数\textbf{发散}。
\end{enumerate}

综上结论：
$p>1$ 或 $p=1$ 且 $q>1$ 时，级数收敛；
其余情况级数发散。
\end{example}

\begin{example}{}{}
讨论级数 $\sum_{n=1}^\infty (-1)^n \cdot \frac{\arctan n}{n^p}$ 的敛散性。
\tcblower
通项为交错符号+有界因子+幂次分母，考虑Abel判别法——若级数 $\sum a_n$ 收敛，数列 $\{b_n\}$ 单调有界，则级数 $\sum a_n b_n$ 收敛。先判断绝对收敛性，再用Abel判别法判断条件收敛性。\\
设 $a_n = \frac{(-1)^n}{n^p}$，$b_n = \arctan n$。
显然数列 $\{b_n\}$ 单调递增，且 $\lim_{n\to\infty} b_n = \frac{\pi}{2}$，故 $\{b_n\}$ 单调有界。

\begin{enumerate}
    \item \textbf{绝对收敛性分析：}
    正项级数 $\sum \left| (-1)^n \cdot \frac{\arctan n}{n^p} \right| = \sum \frac{\arctan n}{n^p}$。
    当 $n\to\infty$ 时，$\arctan n \to \frac{\pi}{2}$，故 $\frac{\arctan n}{n^p} \sim \frac{\pi}{2 n^p}$。
    \begin{itemize}
        \item 若 $p>1$：$\sum \frac{1}{n^p}$ 收敛，由比较判别法，原级数\textbf{绝对收敛}；
        \item 若 $p\leq1$：$\sum \frac{1}{n^p}$ 发散，故原级数不绝对收敛。
    \end{itemize}

    \item \textbf{条件收敛性分析（$0 < p \leq 1$）：}
    交错级数 $\sum a_n = \sum \frac{(-1)^n}{n^p}$，由莱布尼茨判别法，当 $0 < p \leq 1$ 时收敛。
    又 $\{b_n\}$ 单调有界，由\textbf{Abel判别法}，级数 $\sum a_n b_n$ 收敛。
    结合不绝对收敛的结论，此时级数\textbf{条件收敛}。

    \item \textbf{发散性分析（$p \leq 0$）：}
    通项的绝对值 $\frac{\arctan n}{n^p} \geq \arctan 1 \cdot n^{-p} \geq \frac{\pi}{4}$，不趋于0，由级数收敛的必要条件，原级数\textbf{发散}。
\end{enumerate}

综上结论：
 $p > 1$：级数绝对收敛；
 $0 < p \leq 1$：级数条件收敛；
 $p \leq 0$：级数发散。
\end{example}


\section{高频易错点}
\begin{mistake}{把 $u_n\to0$ 当作收敛判别}{mistake1}
$u_n\to0$ 只是必要条件，不是充分条件。例如调和级数 $\sum 1/n$ 的通项趋于 $0$，但部分和发散。考试中若只写“通项趋零，所以收敛”，属于原则性错误。
\end{mistake}
\begin{mistake}{比值判别临界值为 $1$ 时乱下结论}{mistake2}
若 $\lim |u_{n+1}/u_n|=1$，比值判别失效。例如 $\sum 1/n$ 与 $\sum 1/n^2$ 的比值极限都为 $1$，但一个发散一个收敛。
\end{mistake}
\coursekey\ 课件第12章的核心主线是“级数收敛性”，所有技巧都服务于一个目标：控制部分和或尾和。

\chapter{函数项级数与一致收敛}
\section{核心定义}
\begin{definition}{逐点收敛与一致收敛}{uniformdef}
设 $f_n$ 与 $f$ 均定义在集合 $E$ 上。若对每个固定 $x\in E$，有 $f_n(x)\to f(x)$，则称 $f_n$ 在 $E$ 上逐点收敛于 $f$。若对任意 $\varepsilon>0$，存在仅与\(\varepsilon\)有关的正整数$N(\varepsilon)$，使得当 $n\ge N(\varepsilon)$ 时对一切 $x\in E$ 同时有
\[
|f_n(x)-f(x)|<\varepsilon,
\]
则称 $f_n$ 在 $E$ 上一致收敛于 $f$，记作 $f_n\rightrightarrows f$。
\end{definition} 

\begin{definition}{函数项级数一致收敛}{funseries}
函数项级数 $\sum_{n=1}^{\infty}u_n(x)$ 的部分和函数为
\[
S_N(x)=\sum_{n=1}^{N}u_n(x).
\]
若 $S_N(x)$ 在 $E$ 上一致收敛于 $S(x)$，则称该函数项级数在 $E$ 上一致收敛，称 $S$ 为和函数。
\end{definition}

\begin{warningbox}{一致收敛的必要性}
    注意，仅有点态收敛我们无法得到更好的性质，例如交换求极限的顺序(这里显然是面对无限个函数的情况)：\[
    \lim_{x\to x_0}[u_1(x)+u_2(x)+\cdots u_{\infty}(x)]=\sum_{k=1}^{\infty}\lim_{x\to x_0}u_k(x)
    \]交换积分的顺序：\[
    \int[\sum_{k=1}^{\infty}u_k(x)]\dd x=\sum_{k=1}^{\infty}\int u_k(x)\dd x
    \]交换求导的顺序：\[
    \frac{d}{\dd x}[\sum_{k=1}^{\infty}u_k(x)]=\sum_{k=1}^{\infty}\frac{\dd u_k(x)}{\dd x}
    \]可以举出三种情况的反例，这里不再赘述，故我们需要更强的条件(也就是一致收敛)来达到这样的效果。
\end{warningbox}




\section{关键定理与证明脉络}

\begin{theorem}{一致收敛的充要条件}{uniform}
    函数列\(S_n(x)\)在集合D上点态收敛到\(S(x)\)，那么\(S_n(x)\)一致收敛到\(S(x)\)的充要条件是\[
    \lim_{n\to \infty}\sup|S_n(x)-S(x)|=0
    \]
\end{theorem}



\begin{theorem}{一致收敛的Cauchy准则}{uniformcauchy}
函数项级数 $\sum u_n(x)$ 在 $E$ 上一致收敛的充要条件为：任给 $\varepsilon>0$，存在 $N$，使得当 $n>m\ge N$ 时，
\[
\sup_{x\in E}\left|\sum_{k=m+1}^{n}u_k(x)\right|<\varepsilon.
\]
\end{theorem}

\begin{theorem}{Weierstrass判别法}{weier}
若存在正项级数 $\sum M_n$ 收敛，且对所有 $x\in E$ 有 $|u_n(x)|\le M_n$，则 $\sum u_n(x)$ 在 $E$ 上绝对一致收敛。
\end{theorem}
\noindent\textbf{证明思路：} 对任意尾和，
\[
\left|\sum_{k=m+1}^{n}u_k(x)\right|\le \sum_{k=m+1}^{n}|u_k(x)|\le \sum_{k=m+1}^{n}M_k.
\]
右端与 $x$ 无关并可由正项级数的尾和趋零控制，因此得到一致Cauchy准则。

\begin{theorem}{abel判别法、dirichlet判别法}{abel-dirichlet}
\noindent \textbf{abel判别法：}设函数项级数 $\sum_{n=1}^{\infty} a_n(x) b_n(x)$ 在集合 $E$ 上满足以下条件：

\begin{enumerate}
    \item 级数 $\sum_{n=1}^{\infty} a_n(x)$ 在 $E$ 上一致收敛；
    \item 对每个 $x \in E$，序列 $\{b_n(x)\}$ 关于 $n$ 单调，且在 $E$ 上一致有界，即存在常数 $M > 0$，使得
    \[
    |b_n(x)| \le M, \quad \forall n \in \mathbb{N}, \forall x \in E.
    \]
\end{enumerate}

则级数 $\sum_{n=1}^{\infty} a_n(x) b_n(x)$ 在 $E$ 上一致收敛。\par
\textbf{dirichlet判别法：}设函数项级数 $\sum_{n=1}^{\infty} a_n(x) b_n(x)$ 在集合 $E$ 上满足以下条件：

\begin{enumerate}
    \item 级数 $\sum_{n=1}^{\infty} a_n(x)$ 的部分和序列在 $E$ 上一致有界，即存在常数 $M > 0$，使得
    \[
    \left| \sum_{k=1}^{n} a_k(x) \right| \le M, \quad \forall n \in \mathbb{N}, \forall x \in E;
    \]
    \item 对每个 $x \in E$，序列 $\{b_n(x)\}$ 关于 $n$ 单调，且在 $E$ 上一致收敛于 0。
\end{enumerate}

则级数 $\sum_{n=1}^{\infty} a_n(x) b_n(x)$ 在 $E$ 上一致收敛。


\end{theorem}

\begin{theorem}{一致收敛下的极限、积分、求导交换}{interchange}
设 $f_n\rightrightarrows f$。
\begin{enumerate}
\item 若每个 $f_n$ 在点 $x_0$ 附近连续，那么\(f\)在\(x_0\)附近连续，可以交换$n\to\infty$ 与 $x\to x_0$ 的极限。
\item 若 $f_n\in C[a,b]$ 则
\[
\int_a^b f(x)\dd x=\lim_{n\to\infty}\int_a^b f_n(x)\dd x=\int_{a}^{b}\lim_{n\to\infty}f_n(x)\dd x 
\]
\item 若 每个\(f_n(x)\)都有连续导函数且$\sum f_n'(x)$ 在 $[a,b]$ 上一致收敛，则其和函数可逐项求导。
\end{enumerate}
\end{theorem}

\begin{theorem}{幂级数的基本性质(abel定理)}{powerseries}
幂级数 $\sum a_n(x-x_0)^n$ 存在收敛半径 $R$。在 $|x-x_0|<R$ 内绝对收敛，在任意内闭区间 $|x-x_0|\le r<R$ 上一致收敛，并可逐项求导、逐项积分。端点 $x=x_0\pm R$ 必须代回原级数单独讨论。
\end{theorem}

\section{典型解题方法}
\begin{methodbox}{证明一致收敛的常用路线}
\begin{enumerate}
\item \textbf{上确界路线：} 直接估计 $\sup_{x\in E}|f_n(x)-f(x)|$ 或尾和上确界。
\item \textbf{M判别路线：} 找与 $x$ 无关的 $M_n$，使 $|u_n(x)|\le M_n$，再证明 $\sum M_n$ 收敛。
\item \textbf{Dirichlet一致路线：} 证明 $\sum_{k=1}^{n}a_k(x)$ 的部分和关于 $x$ 一致有界，且 $b_n(x)$ 关于 $n$ 单调并一致趋零。
\item \textbf{反证构造路线：} 若要证明不一致，通常找 $x_n$，使 $|f_n(x_n)-f(x_n)|\not\to0$。
\end{enumerate}
\end{methodbox}

\begin{example}{典型例：$x^n$ 在不同区间上的收敛性}{xn}
函数列 $f_n(x)=x^n$ 在 $[0,1]$ 上逐点收敛到
\[
f(x)=\begin{cases}0,&0\le x<1,\\1,&x=1.
\end{cases}
\]
极限函数不连续，而每个 $x^n$ 连续，所以不可能一致收敛。在 $[0,a]$，$0<a<1$ 上，
\[
\sup_{0\le x\le a}|x^n|=a^n\to0,
\]
故一致收敛。
\end{example}

\begin{example}{点态收敛域求解}{solve}
求下列函数项级数的收敛域：
\begin{enumerate}
  \item $\displaystyle \sum_{n=1}^{\infty} \frac{(x-1)^n}{n \cdot 2^n}$
  \item $\displaystyle \sum_{n=1}^{\infty} \frac{1}{n} \left( \frac{x-1}{x+1} \right)^n$
  \item $\displaystyle \sum_{n=1}^{\infty} \frac{\cos(nx)}{n^x}$
\end{enumerate}

\smallskip
\noindent \textbf{考点}：函数项级数点态收敛的基本思路、幂级数收敛域、狄利克雷判别法的应用。

(1) 该级数为幂级数，由比值判别法求收敛半径：
$$
\lim_{n \to \infty} \left| \frac{a_{n+1}}{a_n} \right| = \lim_{n \to \infty} \frac{n \cdot 2^n}{(n+1) \cdot 2^{n+1}} = \frac{1}{2}
$$
故收敛半径 $R=2$，收敛区间为 $|x-1|<2$，即 $x \in (-1,3)$。

- 端点 $x=3$：级数为 $\sum_{n=1}^\infty \frac{1}{n}$（调和级数），发散；
- 端点 $x=-1$：级数为 $\sum_{n=1}^\infty \frac{(-1)^n}{n}$（莱布尼茨级数），收敛。

因此收敛域为 $\boldsymbol{[-1, 3)}$。

\medskip
(2) 令 $t = \frac{x-1}{x+1}$，级数化为 $\sum_{n=1}^\infty \frac{t^n}{n}$，其收敛域为 $t \in [-1,1)$。

解不等式 $-1 \leq \frac{x-1}{x+1} < 1$：
- 右侧：$\frac{x-1}{x+1} < 1 \implies \frac{-2}{x+1} < 0 \implies x > -1$；
- 左侧：$\frac{x-1}{x+1} \geq -1 \implies \frac{2x}{x+1} \geq 0 \implies x \geq 0 \text{ 或 } x < -1$。

取交集得 $x \geq 0$，故收敛域为 $\boldsymbol{[0, +\infty)}$。

\medskip
(3) 分情况讨论：
- 当 $x > 1$ 时：$\left| \frac{\cos(nx)}{n^x} \right| \leq \frac{1}{n^x}$，p级数 $\sum \frac{1}{n^x}$ 收敛，原级数绝对收敛；
- 当 $0 < x \leq 1$ 时：对 $x \neq 2k\pi\ (k \in \mathbb{Z})$，部分和 $\left| \sum_{k=1}^n \cos(kx) \right| \leq \frac{1}{|\sin(x/2)|}$ 一致有界，数列 $\{\frac{1}{n^x}\}$ 单调趋于0，由狄利克雷判别法，级数收敛；
  对 $x = 2k\pi$，级数化为 $\sum \frac{1}{n^x}$，$0<x\leq1$ 时发散；
- 当 $x \leq 0$ 时：通项 $\frac{\cos(nx)}{n^x}$ 不趋于0，级数发散。

综上，收敛域为 $\boldsymbol{(0, +\infty)}$（$x=2k\pi$ 处发散，其余点收敛）。


\end{example}

\begin{example}{魏尔斯特拉斯M判别法}{weierstrass}
\noindent 证明下列级数在指定区间上一致收敛：
\begin{enumerate}
  \item $\displaystyle \sum_{n=1}^{\infty} \frac{\sin(nx) + \cos(n^2 x)}{n^{3/2}},\quad x \in \mathbb{R}$
  \item $\displaystyle \sum_{n=1}^{\infty} x^2 e^{-nx},\quad x \in [0, +\infty)$
\end{enumerate}
\tcblower
\smallskip
\noindent \textbf{考点}：M判别法的使用条件、函数最值求解、绝对一致收敛的判定。


M判别法核心：若对所有 $x \in I$，有 $|u_n(x)| \leq M_n$，且数项级数 $\sum M_n$ 收敛，则 $\sum u_n(x)$ 在 $I$ 上一致收敛。

\medskip
(1) 对任意 $x \in \mathbb{R}$，由三角函数有界性：
$$
\left| \frac{\sin(nx) + \cos(n^2 x)}{n^{3/2}} \right| \leq \frac{|\sin(nx)| + |\cos(n^2 x)|}{n^{3/2}} \leq \frac{2}{n^{3/2}}
$$
p级数 $\sum_{n=1}^\infty \frac{2}{n^{3/2}}$ 收敛（$p=3/2>1$），由M判别法，原级数在 $\mathbb{R}$ 上一致收敛。

\medskip
(2) 先求通项 $u_n(x) = x^2 e^{-nx}$ 在 $[0,+\infty)$ 上的最大值：
求导得
$$
u_n'(x) = 2x e^{-nx} - n x^2 e^{-nx} = x e^{-nx}(2 - nx)
$$
令 $u_n'(x)=0$，得驻点 $x=0$ 与 $x=2/n$。代入得
$$
u_n\left(\frac{2}{n}\right) = \frac{4}{n^2} e^{-2} = \frac{4}{e^2 n^2}
$$
为区间上的最大值。

因此对所有 $x \in [0,+\infty)$，有 $0 \leq u_n(x) \leq \frac{4}{e^2 n^2}$。
而 $\sum_{n=1}^\infty \frac{4}{e^2 n^2}$ 收敛，由M判别法，原级数在 $[0,+\infty)$ 上一致收敛。
\end{example}

\begin{example}{不一致收敛的判定}{a}
\noindent 讨论级数 $\displaystyle \sum_{n=1}^{\infty} x e^{-nx}$ 在下列区间上的一致收敛性：
\begin{enumerate}
  \item $[a, +\infty)$，其中 $a > 0$；
  \item $(0, +\infty)$。
\end{enumerate}
\tcblower
\smallskip
\noindent \textbf{考点}：一致收敛的余项准则、不一致收敛的证明技巧、内闭一致收敛概念。

一致收敛的余项准则：$\sum u_n(x)$ 在 $I$ 上一致收敛 $\iff$ $\displaystyle \lim_{n \to \infty} \sup_{x \in I} |R_n(x)| = 0$，其中 $R_n(x) = \sum_{k=n+1}^\infty u_k(x)$ 为余项。

\medskip
(1) 在 $[a, +\infty)\ (a>0)$ 上：
对通项 $u_n(x) = x e^{-nx}$ 求导，得 $u_n'(x) = e^{-nx}(1 - nx)$，驻点为 $x=1/n$。

当 $n > 1/a$ 时，$1/n < a$，此时 $u_n(x)$ 在 $[a, +\infty)$ 上单调递减，最大值在 $x=a$ 处：
$$
u_n(a) = a e^{-n a}
$$
因此对所有 $x \geq a$，有 $0 < u_n(x) \leq a e^{-n a}$。

等比级数 $\sum_{n=1}^\infty a e^{-n a}$ 收敛（公比 $e^{-a} < 1$），由M判别法，级数在 $[a, +\infty)$ 上一致收敛。

\medskip
(2) 在 $(0, +\infty)$ 上：
先求余项，该级数为等比级数，首项 $x e^{-(n+1)x}$，公比 $e^{-x}$，故：
$$
R_n(x) = \frac{x e^{-(n+1)x}}{1 - e^{-x}} = \frac{x e^{-n x}}{e^x - 1}
$$

取特殊点列 $x = 1/n \in (0,+\infty)$，代入得：
$$
R_n\left( \frac{1}{n} \right) = \frac{\frac{1}{n} \cdot e^{-1}}{e^{1/n} - 1}
$$

当 $n \to \infty$ 时，$e^{1/n} - 1 \sim \frac{1}{n}$，因此：
$$
\lim_{n \to \infty} R_n\left( \frac{1}{n} \right) = e^{-1} \neq 0
$$

这说明 $\displaystyle \sup_{x \in (0,+\infty)} |R_n(x)| \geq e^{-1}$，不趋于0。由余项准则，级数在 $(0, +\infty)$ 上不一致收敛。
\end{example}

\begin{example}{含参级数一致收敛性讨论}{b}
讨论级数 $\displaystyle \sum_{n=1}^{\infty} \frac{x}{n^\alpha (1 + n x^2)}$ 在 $\mathbb{R}$ 上的一致收敛性，其中 $\alpha>0$ 为参数。
\tcblower
\smallskip
\noindent \textbf{考点}：参数分类思想、函数最值分析、余项下界估计。


级数为奇函数，只需讨论 $x \geq 0$ 的情形。

对固定的 $n$ 和 $\alpha>0$，求 $f_n(x) = \frac{x}{n^\alpha (1 + n x^2)}$ 在 $[0,+\infty)$ 上的最大值：
求导得
$$
f_n'(x) = \frac{1 - n x^2}{n^\alpha (1 + n x^2)^2}
$$
令 $f_n'(x)=0$，得驻点 $x = 1/\sqrt{n}$，最大值为：
$$
f_n\left( \frac{1}{\sqrt{n}} \right) = \frac{1}{2 n^{\alpha + 1/2}}
$$

分情况讨论：

1. **当 $\boldsymbol{\alpha > 1/2}$ 时**：
对所有 $x \in \mathbb{R}$，$|u_n(x)| \leq \frac{1}{2 n^{\alpha + 1/2}}$，而p级数 $\sum \frac{1}{n^{\alpha + 1/2}}$ 收敛（$p=\alpha+1/2>1$），由M判别法，级数在 $\mathbb{R}$ 上一致收敛。

2. **当 $\boldsymbol{0 < \alpha \leq 1/2}$ 时**：
取特殊点 $x = 1/\sqrt{n}$，估计余项下界：
$$
R_n\left( \frac{1}{\sqrt{n}} \right) = \sum_{k=n+1}^\infty \frac{1/\sqrt{n}}{k^\alpha (1 + k/n)}
$$

当 $k > n$ 时，$1 + k/n < 2k/n$，因此：
$$
\frac{1}{k^\alpha (1 + k/n)} > \frac{n}{2 k^{\alpha + 1}}
$$

代入余项得：
$$
R_n\left( \frac{1}{\sqrt{n}} \right) > \frac{\sqrt{n}}{2} \sum_{k=n+1}^\infty \frac{1}{k^{\alpha + 1}}
$$

由积分估计，$\sum_{k=n+1}^\infty \frac{1}{k^{\alpha+1}} \sim \frac{1}{\alpha n^\alpha}\ (n \to \infty)$，因此：
$$
R_n\left( \frac{1}{\sqrt{n}} \right) > \frac{1}{2\alpha} n^{\frac{1}{2} - \alpha}
$$

- 若 $0<\alpha<1/2$，则 $n^{\frac{1}{2}-\alpha} \to +\infty$，余项不趋于0；
- 若 $\alpha=1/2$，则 $R_n(1/\sqrt{n}) > 1$，余项不趋于0。

因此 $0<\alpha\leq1/2$ 时，级数在 $\mathbb{R}$ 上不一致收敛。

\medskip
综上：当 $\alpha > 1/2$ 时级数一致收敛；当 $0 < \alpha \leq 1/2$ 时不一致收敛。

\end{example}








\section{高频易错点}
\begin{mistake}{逐点收敛与一致收敛混淆}{uniformmistake}
逐点收敛的 $N$ 可以依赖于 $x$，一致收敛的 $N$ 不能依赖于 $x$。凡是题目问“能否换极限、换积分、换求导”，必须检查一致性，而不能只检查逐点收敛。
\end{mistake}
\begin{mistake}{幂级数端点遗漏}{endpoint}
用比值法或根值法求得收敛半径后，只能确定开区间内外的情况，端点上可能收敛、发散、条件收敛或绝对收敛，需要逐点代入。
\end{mistake}
\examkey\ 2025回忆卷中递推积分函数列题，本质就是用阶乘型估计构造一致收敛的控制级数。

\chapter{重积分}
\section{核心定义}
\begin{definition}{二重积分与三重积分}{multipleint}
二重积分 $\iint_D f(x,y)\dd x\dd y$ 是将区域 $D$ 分割成小块后取积分和极限得到的总量；三重积分 $\iiint_V f(x,y,z)\dd x\dd y\dd z$ 是空间区域上的类似总量。若 $f\equiv1$，二重积分表示面积，三重积分表示体积。
\end{definition}

\section{关键定理与计算公式}
\begin{theorem}{Fubini定理}{fubini}
若 $D$ 是 $x$ 型区域
\[
D=\{(x,y):a\le x\le b,\ \varphi_1(x)\le y\le \varphi_2(x)\},
\]
且 $f$ 连续，则
\[
\iint_D f(x,y)\dd x\dd y=\int_a^b\dd x\int_{\varphi_1(x)}^{\varphi_2(x)}f(x,y)\dd y.
\]
$y$ 型区域有类似公式。三重积分可按投影区域与上下曲面写成累次积分。
\end{theorem}

\begin{theorem}{变量替换公式}{changevar}
若 $(x,y)=\Phi(u,v)$ 在区域 $D'$ 上一一、连续可微，Jacob行列式
\[
J=\frac{\partial(x,y)}{\partial(u,v)}
\]
不为零，则
\[
\iint_D f(x,y)\dd x\dd y=\iint_{D'} f(x(u,v),y(u,v))|J|\dd u\dd v.
\]
三维变量替换同理。(可以认为这里的\(|J|\)是换元之后的面积变换比例)
\end{theorem}

\begin{formula}{常用坐标公式}
\begin{align*}
\text{极坐标：}\quad &x=r\cos\theta,\quad y=r\sin\theta,
&\dd x\dd y&=r\dd r\dd\theta.\\
\text{柱坐标：}\quad &x=r\cos\theta,
\quad y=r\sin\theta,
\quad z=z,
&\dd V&=r\dd r\dd\theta\dd z.\\
\text{球坐标：}\quad &x=\rho\sin\varphi\cos\theta,
\quad y=\rho\sin\varphi\sin\theta,
\quad z=\rho\cos\varphi,
&\dd V&=\rho^2\sin\varphi\dd\rho\dd\varphi\dd\theta.
\end{align*}
\end{formula}

\section{典型解题方法}
\begin{methodbox}{重积分区域处理模板}
\begin{enumerate}
\item \textbf{先画区域或投影。} 不画图直接写积分限，是重积分错误的主要来源。
\item \textbf{看对称性。} 区域关于某轴或某面对称，被积函数为奇函数时对应积分为零。
\item \textbf{看边界类型。} 圆、扇形、球、圆柱优先用极/柱/球坐标；抛物线、椭圆可考虑特殊换元。
\item \textbf{检查Jacobian。} 每次换元都必须写出Jacobian，且注意取绝对值。
\item \textbf{检查积分限方向。} 内层变量上下限必须与外层变量取值一致，不得出现空区间。
\end{enumerate}
\end{methodbox}

\begin{example}{典型例：球与圆柱交叠体积}{cylinderball}
求球 $x^2+y^2+z^2\le a^2$ 被圆柱 $x^2+y^2\le ax$ 截得部分的体积。柱坐标下，圆柱为
\[
r^2\le ar\cos\theta\quad\Longleftrightarrow\quad 0\le r\le a\cos\theta,
\]
其中 $-\pi/2\le\theta\le\pi/2$。球给出 $-\sqrt{a^2-r^2}\le z\le\sqrt{a^2-r^2}$，故
\[
V=\int_{-\pi/2}^{\pi/2}\int_0^{a\cos\theta}2\sqrt{a^2-r^2}\,r\dd r\dd\theta.
\]
若继续计算，先对 $r$ 积分：
\[
\int 2r\sqrt{a^2-r^2}\dd r=-\frac{2}{3}(a^2-r^2)^{3/2}.
\]
\end{example}

\begin{example}{二重积分交换积分次序}{c}
交换下列二重积分的积分次序：
\begin{enumerate}
  \item $\displaystyle \int_{0}^{1} dx \int_{x^2}^{x} f(x,y) dy$
  \item $\displaystyle \int_{0}^{2} dx \int_{0}^{\frac{x^2}{2}} f(x,y) dy + \int_{2}^{2\sqrt{2}} dx \int_{0}^{\sqrt{8-x^2}} f(x,y) dy$
\end{enumerate}
\tcblower
\smallskip
\noindent \textbf{考点}：积分区域的几何表示、X型区域与Y型区域的转换、分段积分区域的合并。

交换积分次序的核心步骤：根据原积分限画出积分区域，再按另一种坐标型重新确定积分限。

\medskip
(1) 原积分是X型区域：$0\leq x\leq 1,\ x^2\leq y\leq x$，区域由抛物线$y=x^2$与直线$y=x$围成。

转换为Y型区域：$0\leq y\leq 1,\ y\leq x\leq \sqrt{y}$，因此交换后积分：
$$
\int_{0}^{1} dy \int_{y}^{\sqrt{y}} f(x,y) dx
$$

\medskip
(2) 原积分由两个X型区域组成：
- 区域$D_1$：$0\leq x\leq 2,\ 0\leq y\leq \frac{x^2}{2}$（抛物线下方）
- 区域$D_2$：$2\leq x\leq 2\sqrt{2},\ 0\leq y\leq \sqrt{8-x^2}$（四分之一圆下方）

合并为整体区域，转换为Y型区域：$0\leq y\leq 2,\ \sqrt{2y}\leq x\leq \sqrt{8-y^2}$，因此交换后积分：
$$
\int_{0}^{2} dy \int_{\sqrt{2y}}^{\sqrt{8-y^2}} f(x,y) dx
$$
\end{example}

\begin{example}{二重积分的变量替换}{d}
计算下列二重积分：
\begin{enumerate}
  \item $\displaystyle \iint_D e^{\frac{x-y}{x+y}} \, dxdy$，其中$D$由$x=0,\ y=0,\ x+y=1$围成的三角形区域
  \item $\displaystyle \iint_D \sqrt{1-\frac{x^2}{a^2}-\frac{y^2}{b^2}} \, dxdy$，其中$D$为椭圆域 $\frac{x^2}{a^2}+\frac{y^2}{b^2} \leq 1$
\end{enumerate}
\tcblower
\smallskip
\noindent \textbf{考点}：一般变量替换公式、雅可比行列式的计算、广义极坐标变换。

二重积分变量替换公式：设变换$x=x(u,v),\ y=y(u,v)$，雅可比行列式$J = \frac{\partial(x,y)}{\partial(u,v)} \neq 0$，则
$$
\iint_D f(x,y) dxdy = \iint_{D'} f[x(u,v),y(u,v)] \cdot |J| \, dudv
$$

\medskip
(1) 令$u = x+y,\ v = x-y$，解得$x = \frac{u+v}{2},\ y = \frac{u-v}{2}$。

计算雅可比行列式：
$$
J = \frac{\partial(x,y)}{\partial(u,v)} = \begin{vmatrix} \frac{1}{2} & \frac{1}{2} \\ \frac{1}{2} & -\frac{1}{2} \end{vmatrix} = -\frac{1}{2}, \quad |J| = \frac{1}{2}
$$

原区域$D$对应$uOv$平面区域$D'$：$0\leq u\leq 1,\ -u\leq v\leq u$。因此：
$$
\begin{aligned}
\iint_D e^{\frac{x-y}{x+y}} dxdy &= \iint_{D'} e^{\frac{v}{u}} \cdot \frac{1}{2} dudv \\
&= \frac{1}{2} \int_{0}^{1} du \int_{-u}^{u} e^{\frac{v}{u}} dv \\
&= \frac{1}{2} \int_{0}^{1} u \cdot \left( e - e^{-1} \right) du \\
&= \frac{e - e^{-1}}{2} \cdot \frac{1}{2} = \boldsymbol{\frac{e - e^{-1}}{4}}
\end{aligned}
$$

\medskip
(2) 采用广义极坐标变换：$x = ar\cos\theta,\ y = br\sin\theta$。

计算雅可比行列式：
$$
J = \frac{\partial(x,y)}{\partial(r,\theta)} = \begin{vmatrix} a\cos\theta & -ar\sin\theta \\ b\sin\theta & br\cos\theta \end{vmatrix} = abr, \quad |J| = abr
$$

椭圆域对应$0\leq \theta\leq 2\pi,\ 0\leq r\leq 1$，因此：
$$
\begin{aligned}
\iint_D \sqrt{1-\frac{x^2}{a^2}-\frac{y^2}{b^2}} dxdy &= \int_{0}^{2\pi} d\theta \int_{0}^{1} \sqrt{1-r^2} \cdot abr \, dr \\
&= 2\pi ab \cdot \left. \left( -\frac{1}{3}(1-r^2)^{\frac{3}{2}} \right) \right|_{0}^{1} \\
&= \boldsymbol{\frac{2\pi ab}{3}}
\end{aligned}
$$
\end{example}
\begin{example}{重积分的几何与物理应用}
设均匀物体（密度$\rho=1$）由曲面$z=x^2+y^2$与$z=2-\sqrt{x^2+y^2}$围成，求：
\begin{enumerate}
  \item 该物体的体积；
  \item 该物体的质心坐标；
  \item 曲面$z=x^2+y^2$被截部分的曲面面积。
\end{enumerate}
\tcblower
\smallskip
\noindent \textbf{考点}：重积分求体积、曲面面积公式、质心坐标公式、综合应用能力。

先求两曲面的交线：联立$z=x^2+y^2$与$z=2-\sqrt{x^2+y^2}$，令$r=\sqrt{x^2+y^2}$，得$r^2 + r - 2 = 0$，解得$r=1$，即交线为$z=1$平面上的圆$x^2+y^2=1$。

\medskip
(1) 体积计算（柱坐标）：
区域$\Omega$：$0\leq \theta\leq 2\pi,\ 0\leq r\leq 1,\ r^2\leq z\leq 2-r$
$$
\begin{aligned}
V &= \iiint_\Omega dV = \int_{0}^{2\pi} d\theta \int_{0}^{1} r dr \int_{r^2}^{2-r} dz \\
&= 2\pi \int_{0}^{1} r(2-r-r^2) dr \\
&= 2\pi \int_{0}^{1} (2r - r^2 - r^3) dr \\
&= 2\pi \cdot \left( 1 - \frac{1}{3} - \frac{1}{4} \right) = \boldsymbol{\frac{5\pi}{6}}
\end{aligned}
$$

\medskip
(2) 质心坐标：
由对称性，物体关于$z$轴对称，故质心的$x,y$坐标$\bar{x}=0,\ \bar{y}=0$，只需计算$\bar{z} = \frac{1}{V} \iiint_\Omega z dV$。

计算静矩：
$$
\begin{aligned}
\iiint_\Omega z dV &= \int_{0}^{2\pi} d\theta \int_{0}^{1} r dr \int_{r^2}^{2-r} z dz \\
&= \pi \int_{0}^{1} r\left[ (2-r)^2 - r^4 \right] dr \\
&= \pi \int_{0}^{1} (4r - 4r^2 + r^3 - r^5) dr \\
&= \pi \cdot \left( 2 - \frac{4}{3} + \frac{1}{4} - \frac{1}{6} \right) = \frac{5\pi}{4}
\end{aligned}
$$

因此：
$$
\bar{z} = \frac{\frac{5\pi}{4}}{\frac{5\pi}{6}} = \frac{3}{2}
$$
质心坐标为$\boldsymbol{\left( 0, 0, \frac{3}{2} \right)}$。

\medskip
(3) 曲面面积计算：
曲面$z=x^2+y^2$的面积元素$dS = \sqrt{1+\left( \frac{\partial z}{\partial x} \right)^2 + \left( \frac{\partial z}{\partial y} \right)^2} dxdy = \sqrt{1+4x^2+4y^2} dxdy$。

投影区域$D$：$x^2+y^2 \leq 1$，极坐标下计算：
$$
\begin{aligned}
S &= \iint_D \sqrt{1+4(x^2+y^2)} dxdy \\
&= \int_{0}^{2\pi} d\theta \int_{0}^{1} \sqrt{1+4r^2} \cdot r dr \\
&= 2\pi \cdot \frac{1}{12} \left. (1+4r^2)^{\frac{3}{2}} \right|_{0}^{1} \\
&= \frac{\pi}{6} \left( 5\sqrt{5} - 1 \right)
\end{aligned}
$$
即曲面面积为$\boldsymbol{\frac{\pi}{6}\left( 5\sqrt{5} - 1 \right)}$。
\end{example}







\section{高频易错点}
\begin{mistake}{极角范围与半径上限不匹配}{polarmistake}
例如 $x^2+y^2\le ax$ 在极坐标中是 $0\le r\le a\cos\theta$，因此必须有 $\cos\theta\ge0$，即 $-\pi/2\le\theta\le\pi/2$。若写成 $0\le\theta\le2\pi$ 会导致半径上限为负。
\end{mistake}
\begin{mistake}{换元后漏掉Jacobian}{jacmistake}
变量替换不是单纯替换被积函数，面积元、体积元也要变。考试中“式子看似正确但差一个 $r$ 或 $\rho^2\sin\varphi$”是最常见失分点。
\end{mistake}

\chapter{曲线积分与曲面积分}
\section{核心定义}
\begin{definition}{第一型与第二型曲线积分}{lineints}
第一型曲线积分 $\int_L f(x,y)\dd s$ 是沿曲线对弧长的积分，与方向无关。若 $L$ 参数化为 $x=x(t),y=y(t)$，则
\[
\int_L f\dd s=\int_\alpha^\beta f(x(t),y(t))\sqrt{x'(t)^2+y'(t)^2}\dd t.
\]
第二型曲线积分
\[
\int_L P\dd x+Q\dd y
\]
与方向有关，参数化后为
\[
\int_\alpha^\beta [P(x(t),y(t))x'(t)+Q(x(t),y(t))y'(t)]\dd t.
\]
\end{definition}
第一类曲线积分的物理意义是曲线的质量，故与方向无关；而第二类曲线积分的物理意义是做功，所以与方向有关。
\begin{definition}{曲面积分与通量}{surfaceints}
设曲面的参数方程为\[\begin{cases}
    x=x(u,v)\\
    y=y(u,v)\\
    z=z(u,v)\\
\end{cases}\]第一型曲面积分 $\iint_S f\dd S$ 与侧向无关。计算方法为\[
    \iint_{\sum} f(x,y,z)\dd S=\iint_{D}f(x(u,v),y(u,v),z(u,v))\sqrt{EG-F^2}\dd u\dd v
    \]其中\[
        \begin{cases}
            E=x_u^2+y_u^2+z_u^2\\
            F=x_ux_v+y_uy_v+z_uz_v\\
            G=x_v^2+y_v^2+z_v^2\\
        \end{cases}
    \]
第二型曲面积分
\[
\iint_{\sigma} P\dd y\dd z+Q\dd z\dd x+R\dd x\dd y
\]
可理解为向量场 $\bm F=(P,Q,R)$ 穿过有向曲面 $S$ 的通量：
\[
\iint_S \bm F\cdot \bm n\dd S.
\]计算方法为\[
 \begin{align*}
    \iint_{\sum} P\dd y\dd z+Q\dd z\dd x+R\dd x\dd y&=\iint_{\sum}[P\cos\alpha+Q\cos\beta+R\cos\gamma]\dd S\\
     &=\pm[P(x(u,v),y(u,v),z(u,v))\frac{\partial(y,z)}{\partial(u,v)}\\
     &+Q(x(u,v),y(u,v),z(u,v))\frac{\partial(x,z)}{\partial(u,v)}\\
     &+R(x(u,v),y(u,v),z(u,v))\frac{\partial(x,y)}{\partial(u,v)}\dd u \dd v]
 \end{align*}\]这里\(\frac{\partial(y,z)}{\partial(u,v)}\)表示\(y,z\)对\(u,v\)的Jacobi行列式，即\[
     \displaystyle\begin{vmatrix}
         \frac{\partial y}{\partial u}&\frac{\partial y}{\partial v}\\
         \frac{\partial z}{\partial u}&\frac{\partial z}{\partial v}\\
     \end{vmatrix}
 \]最前面的正负号由单位法向量与曲面的测决定。
\end{definition}
第一类曲面积分的物理意义是计算曲面面积，而第二类曲面积分的物理意义是计算流量。
\section{关键定理}
\begin{theorem}{Green公式}{green}
设 $D$ 为平面有界区域，边界 $\partial D$ 分段光滑并取正向，$P,Q$ 在 $D$ 上有连续偏导                                                             数，则
\[
\oint_{\partial D}P\dd x+Q\dd y=\iint_D\left(\frac{\partial Q}{\partial x}-\frac{\partial P}{\partial y}\right)\dd x\dd y.
\]
\end{theorem}

\begin{theorem}{Gauss公式}{gauss}
设 $V$ 为空间有界区域，边界 $\partial V$ 取外侧，$\bm F=(P,Q,R)$ 有连续偏导数，则
\[
\iint_{\partial V}\bm F\cdot\bm n\dd S=\iiint_V\diver\bm F\dd V.
\]即\[
\iiint_{V}(\frac{\partial P}{\partial x}+\frac{\partial Q}{\partial y}+\frac{\partial R}{\partial z})\dd x\dd y\dd z=\iint_{\partial V}P\dd y\dd z+Q\dd z\dd x+R\dd x\dd y
\]
\end{theorem}

\begin{theorem}{Stokes公式}{stokes}
设有向曲面 $S$ 的边界为 $\partial S$，方向满足右手法则，则
\[
\oint_{\partial S}\bm F\cdot\dd\bm r=\iint_S(\curl\bm F)\cdot\bm n\dd S.
\]即\[
\iint_{\partial S}=P\dd x+Q\dd y+R\dd z=\iint_{S}(\frac{\partial R}{\partial y}-\frac{\partial Q}{\partial z})\dd y\dd z+(\frac{\partial P}{\partial z}-\frac{\partial R}{\partial x})\dd z\dd x+(\frac{\partial Q}{\partial }-\frac{\partial P}{\partial y})\dd x\dd y
\]
\end{theorem}

\begin{theorem}{调和函数在圆盘上的面积平均值性质}{harmonic-area-mean}
设 $u(x,y)$ 在圆盘
\[
  D_R=\{(x,y):(x-a)^2+(y-b)^2<R^2\}
\]
内二阶连续可微，并且满足
\[
  \Delta u=u_{xx}+u_{yy}=0.
\]
则对任意 $0<\rho<R$，有
\[
  \frac{1}{\pi\rho^2}
  \iint_{(x-a)^2+(y-b)^2\leq \rho^2}u(x,y)\,\mathrm{d}x\mathrm{d}y
  =
  u(a,b).
\]
也就是说，调和函数在圆盘上的面积平均值等于圆心处函数值。
\end{theorem}

\begin{example}{面积平均值性质的详细推导}{ex-harmonic-area-mean-proof}
设 $u(x,y)$ 在圆盘
\[
  D_R=\{(x,y):(x-a)^2+(y-b)^2<R^2\}
\]
内调和，即
\[
  \Delta u=u_{xx}+u_{yy}=0.
\]
证明
\[
  \frac{1}{\pi\rho^2}
  \iint_{(x-a)^2+(y-b)^2\leq \rho^2}u(x,y)\,\mathrm{d}x\mathrm{d}y
  =
  u(a,b),
  \qquad 0<\rho<R.
\]
\tcblower
先证明圆周平均值性质。对 $0<r<R$，记
\[
  M(r)=\frac{1}{2\pi}\int_0^{2\pi}
  u(a+r\cos\theta,b+r\sin\theta)\,\mathrm{d}\theta.
\]
我们证明 $M(r)$ 与 $r$ 无关。对 $r$ 求导，得
\[
  M'(r)
  =
  \frac{1}{2\pi}\int_0^{2\pi}
  \frac{\partial}{\partial r}
  u(a+r\cos\theta,b+r\sin\theta)\,\mathrm{d}\theta.
\]
由链式法则，
\[
  \frac{\partial}{\partial r}
  u(a+r\cos\theta,b+r\sin\theta)
  =
  u_x(a+r\cos\theta,b+r\sin\theta)\cos\theta
  +
  u_y(a+r\cos\theta,b+r\sin\theta)\sin\theta.
\]
注意在圆周
\[
  C_r:(x-a)^2+(y-b)^2=r^2
\]
上，外法向单位向量为
\[
  \mathbf n=(\cos\theta,\sin\theta).
\]
所以
\[
  u_x\cos\theta+u_y\sin\theta
  =
  \frac{\partial u}{\partial n}.
\]
因此
\[
  M'(r)
  =
  \frac{1}{2\pi}\int_0^{2\pi}
  \frac{\partial u}{\partial n}\,\mathrm{d}\theta.
\]
又因为圆周弧长微元满足
\[
  \mathrm{d}s=r\,\mathrm{d}\theta,
\]
所以
\[
  \int_0^{2\pi}\frac{\partial u}{\partial n}\,\mathrm{d}\theta
  =
  \frac{1}{r}\int_{C_r}\frac{\partial u}{\partial n}\,\mathrm{d}s.
\]
于是
\[
  M'(r)
  =
  \frac{1}{2\pi r}\int_{C_r}\frac{\partial u}{\partial n}\,\mathrm{d}s.
\]

由 Green 公式的散度形式，
\[
  \int_{C_r}\frac{\partial u}{\partial n}\,\mathrm{d}s
  =
  \iint_{D_r}\Delta u\,\mathrm{d}x\mathrm{d}y,
\]
其中
\[
  D_r=\{(x,y):(x-a)^2+(y-b)^2\leq r^2\}.
\]
由于 $u$ 是调和函数，
\[
  \Delta u=0.
\]
所以
\[
  \int_{C_r}\frac{\partial u}{\partial n}\,\mathrm{d}s=0.
\]
因此
\[
  M'(r)=0.
\]
这说明 $M(r)$ 是常数。

下面求这个常数。因为 $u$ 在 $(a,b)$ 连续，所以当 $r\to0^+$ 时，
\[
  u(a+r\cos\theta,b+r\sin\theta)\to u(a,b)
\]
且收敛关于 $\theta$ 一致。因此
\[
  \lim_{r\to0^+}M(r)
  =
  \frac{1}{2\pi}\int_0^{2\pi}u(a,b)\,\mathrm{d}\theta
  =
  u(a,b).
\]
由于 $M(r)$ 是常数，故对任意 $0<r<R$，
\[
  M(r)=u(a,b).
\]
即
\[
  \frac{1}{2\pi}\int_0^{2\pi}
  u(a+r\cos\theta,b+r\sin\theta)\,\mathrm{d}\theta
  =
  u(a,b).
\]

现在证明面积平均值性质。对圆盘 $D_\rho$ 作极坐标变换：
\[
  x=a+r\cos\theta,\qquad y=b+r\sin\theta,
\]
其中
\[
  0\leq r\leq \rho,\qquad 0\leq \theta\leq 2\pi.
\]
Jacobi 行列式为
\[
  \left|\frac{\partial(x,y)}{\partial(r,\theta)}\right|=r.
\]
所以
\[
  \iint_{D_\rho}u(x,y)\,\mathrm{d}x\mathrm{d}y
  =
  \int_0^\rho\int_0^{2\pi}
  u(a+r\cos\theta,b+r\sin\theta)\,r\,\mathrm{d}\theta\mathrm{d}r.
\]
把 $r$ 提到外层积分中，得
\[
  \iint_{D_\rho}u(x,y)\,\mathrm{d}x\mathrm{d}y
  =
  \int_0^\rho
  r\left[
  \int_0^{2\pi}
  u(a+r\cos\theta,b+r\sin\theta)\,\mathrm{d}\theta
  \right]\mathrm{d}r.
\]
由刚刚证明的圆周平均值性质，
\[
  \frac{1}{2\pi}
  \int_0^{2\pi}
  u(a+r\cos\theta,b+r\sin\theta)\,\mathrm{d}\theta
  =
  u(a,b).
\]
因此
\[
  \int_0^{2\pi}
  u(a+r\cos\theta,b+r\sin\theta)\,\mathrm{d}\theta
  =
  2\pi u(a,b).
\]
代回面积积分，得到
\[
  \iint_{D_\rho}u(x,y)\,\mathrm{d}x\mathrm{d}y
  =
  \int_0^\rho r\cdot 2\pi u(a,b)\,\mathrm{d}r.
\]
由于 $u(a,b)$ 与 $r$ 无关，
\[
  \iint_{D_\rho}u(x,y)\,\mathrm{d}x\mathrm{d}y
  =
  2\pi u(a,b)\int_0^\rho r\,\mathrm{d}r
  =
  2\pi u(a,b)\cdot \frac{\rho^2}{2}.
\]
所以
\[
  \iint_{D_\rho}u(x,y)\,\mathrm{d}x\mathrm{d}y
  =
  \pi\rho^2u(a,b).
\]
两边除以圆盘面积 $\pi\rho^2$，即得
\[
  \frac{1}{\pi\rho^2}
  \iint_{D_\rho}u(x,y)\,\mathrm{d}x\mathrm{d}y
  =
  u(a,b).
\]

\textbf{思路总结：}
面积平均值性质不是直接硬算出来的，而是先利用 Green 公式证明圆周平均值性质，再把圆盘看成无数个同心圆周的叠加。

\textbf{易错点：}
第一，圆周平均值中的右端是圆心值 $u(a,b)$，不是圆周上的某个特殊点值；第二，面积平均时极坐标里必须带上 Jacobi 因子 $r$；第三，Green 公式中
\[
  \int_{C_r}\frac{\partial u}{\partial n}\,\mathrm{d}s
  =
  \iint_{D_r}\Delta u\,\mathrm{d}x\mathrm{d}y
\]
用的是调和条件 $\Delta u=0$，这是整个证明的核心。
\end{example}

\section{典型解题方法}
\begin{methodbox}{场论计算题标准流程}
\begin{enumerate}
\item \textbf{先判断类型。} 第一型看弧长/面积元，第二型看方向/侧向。
\item \textbf{闭合优先公式。} 平面闭曲线优先Green；闭曲面优先Gauss；空间曲线为曲面边界优先Stokes。
\item \textbf{非闭合要补边/补面。} 补完后用公式，再减去补上的部分。
\item \textbf{奇点要挖洞。} 若场在区域内有奇点，不能直接套Green/Gauss，应去掉小圆/小球或选择避开奇点的辅助边界。
\item \textbf{方向最后核对。} 逆时针为平面正向；闭曲面外侧为正向；补面方向必须与整体外法向一致。
\end{enumerate}
\end{methodbox}

\begin{example}{第一类曲线积分：弧长参数化}{ex-line-first}
设 $L$ 为第一象限内圆弧 $x^2+y^2=a^2$，求
\[
  I=\int_L xy\,\mathrm{d}s.
\]
\tcblower
令
\[
  x=a\cos t,\qquad y=a\sin t,\qquad 0\leq t\leq \frac{\pi}{2}.
\]
则
\[
  \mathrm{d}s=\sqrt{(-a\sin t)^2+(a\cos t)^2}\,\mathrm{d}t=a\,\mathrm{d}t.
\]
所以
\[
  I=\int_0^{\pi/2}a^2\sin t\cos t\cdot a\,\mathrm{d}t
  =a^3\int_0^{\pi/2}\sin t\cos t\,\mathrm{d}t
  =\frac{a^3}{2}.
\]
\textbf{提示：}第一型曲线积分与方向无关，本质是沿曲线对弧长累加。
\end{example}

\begin{example}{第二类曲线积分：常规计算}{reguler}
计算$\displaystyle\int_\Gamma ydx + zdy + xdz$，其中$\Gamma$为从点$A(1,1,1)$到$B(2,3,4)$的空间直线段。

\tcblower
1. 直线参数方程：方向向量$\overrightarrow{AB}=(1,2,3)$，故
$$x=1+t,\ y=1+2t,\ z=1+3t,\quad t\text{从}0\text{到}1$$
2. 求微分：$dx=dt,\ dy=2dt,\ dz=3dt$
3. 代入计算：
$$
\begin{aligned}
\int_\Gamma ydx + zdy + xdz &= \int_0^1 \left[ (1+2t)\cdot1 + (1+3t)\cdot2 + (1+t)\cdot3 \right]dt \\
&= \int_0^1 (6+11t)dt = \frac{23}{2}
\end{aligned}
$$
\end{example}

\begin{example}{含奇点的Green公式：挖洞法}{ex-green-hole}
设 $L$ 为包围原点的正向简单闭曲线，计算
\[
  I=\int_L\frac{-y}{x^2+y^2}\,\mathrm{d}x+\frac{x}{x^2+y^2}\,\mathrm{d}y.
\]
\tcblower
被积函数在原点无定义，不能直接在 $L$ 围成的区域上用 Green 公式。取小圆
\[
  C_\varepsilon:x^2+y^2=\varepsilon^2.
\]
在挖去小圆盘后的区域内，
\[
  \frac{\partial}{\partial x}\frac{x}{x^2+y^2}
  -
  \frac{\partial}{\partial y}\frac{-y}{x^2+y^2}=0.
\]
由 Green 公式可得
\[
  I=\int_{C_\varepsilon^+}\frac{-y}{x^2+y^2}\,\mathrm{d}x+\frac{x}{x^2+y^2}\,\mathrm{d}y.
\]
令
\[
  x=\varepsilon\cos t,\qquad y=\varepsilon\sin t,\qquad 0\leq t\leq 2\pi.
\]
则
\[
  \frac{-y}{x^2+y^2}\,\mathrm{d}x+\frac{x}{x^2+y^2}\,\mathrm{d}y
  =
  \sin^2t\,\mathrm{d}t+\cos^2t\,\mathrm{d}t
  =
  \mathrm{d}t.
\]
所以
\[
  I=\int_0^{2\pi}\mathrm{d}t=2\pi.
\]
\textbf{提示：}不能因为偏导差为 $0$ 就直接说积分为 $0$，原点是奇点。
\end{example}

\begin{example}{第一类曲面积分：圆锥面}{ex-surface-first-cone}
设 $S$ 为圆锥面
\[
  z=\sqrt{x^2+y^2},\qquad 0\leq z\leq h,
\]
求
\[
  I=\iint_S z\,\mathrm{d}S.
\]
\tcblower
取参数
\[
  x=r\cos\theta,\qquad y=r\sin\theta,\qquad z=r,
\]
其中
\[
  0\leq r\leq h,\qquad 0\leq\theta\leq2\pi.
\]
因为
\[
  z_x=\frac{x}{\sqrt{x^2+y^2}},\qquad z_y=\frac{y}{\sqrt{x^2+y^2}},
\]
所以
\[
  \mathrm{d}S=\sqrt{1+z_x^2+z_y^2}\,\mathrm{d}x\mathrm{d}y
  =\sqrt2\,\mathrm{d}x\mathrm{d}y.
\]
又 $\mathrm{d}x\mathrm{d}y=r\,\mathrm{d}r\mathrm{d}\theta$，故
\[
  I=\int_0^{2\pi}\int_0^h r\sqrt2\cdot r\,\mathrm{d}r\mathrm{d}\theta
  =\frac{2\pi\sqrt2}{3}h^3.
\]
\textbf{提示：}第一型曲面积分用 $\mathrm{d}S$，不要误当成平面投影面积。
\end{example}

\begin{example}{第二类曲面积分：投影法}{ex-surface-second-projection}
设 $S$ 为上半球面
\[
  x^2+y^2+z^2=a^2,\qquad z\geq0,
\]
取上侧，求
\[
  I=\iint_S z\,\mathrm{d}x\mathrm{d}y.
\]
\tcblower
由于 $S$ 取上侧，有
\[
  z=\sqrt{a^2-x^2-y^2}.
\]
因此
\[
  I=\iint_D\sqrt{a^2-x^2-y^2}\,\mathrm{d}x\mathrm{d}y,
\]
其中
\[
  D:x^2+y^2\leq a^2.
\]
化为极坐标：
\[
  I=\int_0^{2\pi}\int_0^a \sqrt{a^2-r^2}\,r\,\mathrm{d}r\mathrm{d}\theta.
\]
令 $u=a^2-r^2$，得
\[
  I=2\pi\cdot\frac12\int_0^{a^2}u^{1/2}\,\mathrm{d}u
  =\frac{2\pi a^3}{3}.
\]
\textbf{提示：}取上侧时 $\mathrm{d}x\mathrm{d}y$ 对应正投影；取下侧则要变号。
\end{example}

\begin{example}{第二类曲面积分：硬算}{hard}
计算$\displaystyle\iint_\Sigma z \, dydz + x \, dzdx + y \, dxdy$，其中$\Sigma$是柱面$x^2 + y^2 = 1$被平面$z = 0$和$z = 4$所截部分，方向取外侧。
\tcblower
采用参数方程法直接硬算，步骤如下：

1. 写出柱面$\Sigma$的参数方程
圆柱面$x^2 + y^2 = 1$的标准参数形式为：
$$
\begin{cases}
x = \cos\theta \\
y = \sin\theta \\
z = z
\end{cases}, \quad \theta \in [0, 2\pi],\ z \in [0, 4]
$$
对应的参数平面区域为 $D = \{ (\theta, z) \mid 0 \le \theta \le 2\pi,\ 0 \le z \le 4 \}$。

2. 计算三个雅可比行列式
根据第二类曲面积分的参数计算公式，分别计算三个坐标面投影的雅可比行列式：
$$
\frac{\partial(y,z)}{\partial(\theta,z)} = \begin{vmatrix}
\frac{\partial y}{\partial \theta} & \frac{\partial y}{\partial z} \\
\frac{\partial z}{\partial \theta} & \frac{\partial z}{\partial z}
\end{vmatrix} = \begin{vmatrix}
\cos\theta & 0 \\
0 & 1
\end{vmatrix} = \cos\theta
$$

$$
\frac{\partial(z,x)}{\partial(\theta,z)} = \begin{vmatrix}
\frac{\partial z}{\partial \theta} & \frac{\partial z}{\partial z} \\
\frac{\partial x}{\partial \theta} & \frac{\partial x}{\partial z}
\end{vmatrix} = \begin{vmatrix}
0 & 1 \\
-\sin\theta & 0
\end{vmatrix} = \sin\theta
$$

$$
\frac{\partial(x,y)}{\partial(\theta,z)} = \begin{vmatrix}
\frac{\partial x}{\partial \theta} & \frac{\partial x}{\partial z} \\
\frac{\partial y}{\partial \theta} & \frac{\partial y}{\partial z}
\end{vmatrix} = \begin{vmatrix}
-\sin\theta & 0 \\
\cos\theta & 0
\end{vmatrix} = 0
$$

3. 确定积分符号
柱面外侧的单位法向量为$(x, y, 0) = (\cos\theta, \sin\theta, 0)$，与参数法向量
$\left( \frac{\partial(y,z)}{\partial(\theta,z)},\ \frac{\partial(z,x)}{\partial(\theta,z)},\ \frac{\partial(x,y)}{\partial(\theta,z)} \right)$
方向完全一致，因此积分取正号。

4. 代入公式计算二重积分
将$x = \cos\theta,\ y = \sin\theta,\ z = z$代入被积表达式，转化为参数区域上的二重积分：
$$
\begin{aligned}
\iint_\Sigma z \, dydz + x \, dzdx + y \, dxdy
&= \iint_D \left( z \cdot \cos\theta + \cos\theta \cdot \sin\theta + \sin\theta \cdot 0 \right) d\theta dz \\
&= \int_{0}^{4} dz \int_{0}^{2\pi} \left( z \cos\theta + \cos\theta \sin\theta \right) d\theta
\end{aligned}
$$

分别计算内层两个三角函数积分：
- $\displaystyle \int_{0}^{2\pi} z \cos\theta \, d\theta = z \cdot \sin\theta \bigg|_{0}^{2\pi} = 0$
- $\displaystyle \int_{0}^{2\pi} \cos\theta \sin\theta \, d\theta = \frac{1}{2} \int_{0}^{2\pi} \sin2\theta \, d\theta = -\frac{1}{4} \cos2\theta \bigg|_{0}^{2\pi} = 0$

内层积分整体为$0$，因此最终结果：
$$
\iint_\Sigma z \, dydz + x \, dzdx + y \, dxdy = \int_{0}^{4} 0 \, dz = 0
$$
\end{example}

\begin{example}{第二类曲面积分：较为复杂的计算}{complex}
计算$\displaystyle\iint_\Sigma \frac{1}{x} \, dydz + \frac{1}{y} \, dzdx + \frac{1}{z} \, dxdy$，其中$\Sigma$为椭球面$\frac{x^2}{a^2} + \frac{y^2}{b^2} + \frac{z^2}{c^2} = 1$，方向取外侧。
\tcblower
将积分拆分为三部分，利用对称性分别计算：
$$
I = \iint_\Sigma \frac{1}{x} \, dydz + \frac{1}{y} \, dzdx + \frac{1}{z} \, dxdy = I_1 + I_2 + I_3
$$

\textbf{1. 计算$I_3 = \displaystyle\iint_\Sigma \frac{1}{z} \, dxdy$}
将椭球面按$z$的正负拆分为上下两片：
- $\Sigma_上$：$z = c\sqrt{1 - \frac{x^2}{a^2} - \frac{y^2}{b^2}}$，取上侧，投影符号为正；
- $\Sigma_下$：$z = -c\sqrt{1 - \frac{x^2}{a^2} - \frac{y^2}{b^2}}$，取下侧，投影符号为负。

两片在$xOy$面上的投影区域均为椭圆域：
$$
D_{xy}: \quad \frac{x^2}{a^2} + \frac{y^2}{b^2} \le 1
$$

分别代入投影公式：
$$
\begin{aligned}
\iint_{\Sigma_上} \frac{1}{z} \, dxdy &= \iint_{D_{xy}} \frac{1}{c\sqrt{1 - \frac{x^2}{a^2} - \frac{y^2}{b^2}}} \, dxdy \\
\iint_{\Sigma_下} \frac{1}{z} \, dxdy &= -\iint_{D_{xy}} \frac{1}{-c\sqrt{1 - \frac{x^2}{a^2} - \frac{y^2}{b^2}}} \, dxdy = \iint_{D_{xy}} \frac{1}{c\sqrt{1 - \frac{x^2}{a^2} - \frac{y^2}{b^2}}} \, dxdy
\end{aligned}
$$

两部分相加得：
$$
I_3 = \frac{2}{c} \iint_{D_{xy}} \frac{1}{\sqrt{1 - \frac{x^2}{a^2} - \frac{y^2}{b^2}}} \, dxdy
$$

作广义极坐标变换 $x = ar\cos\theta,\ y = br\sin\theta$，雅可比行列式 $J = abr$，积分区域变为 $r\in[0,1],\ \theta\in[0,2\pi]$，代入得：
$$
\begin{aligned}
\iint_{D_{xy}} \frac{1}{\sqrt{1 - \frac{x^2}{a^2} - \frac{y^2}{b^2}}} \, dxdy
&= \int_{0}^{2\pi} d\theta \int_{0}^{1} \frac{abr}{\sqrt{1 - r^2}} dr \\
&= ab \cdot 2\pi \cdot \left. \left( -\sqrt{1 - r^2} \right) \right|_{0}^{1} \\
&= 2\pi ab
\end{aligned}
$$

因此：
$$
I_3 = \frac{2}{c} \cdot 2\pi ab = \frac{4\pi ab}{c}
$$

\textbf{2. 同理计算其余两个积分}
由椭球面的对称性，完全类似可得：
$$
I_1 = \iint_\Sigma \frac{1}{x} \, dydz = \frac{4\pi bc}{a}, \quad I_2 = \iint_\Sigma \frac{1}{y} \, dzdx = \frac{4\pi ac}{b}
$$

\textbf{3. 合并最终结果}
$$
I = \frac{4\pi bc}{a} + \frac{4\pi ac}{b} + \frac{4\pi ab}{c} = 4\pi \left( \frac{bc}{a} + \frac{ac}{b} + \frac{ab}{c} \right)
$$
\end{example}

\begin{example}{Gauss公式：闭曲面通量}{ex-gauss-sphere}
设 $S$ 为球面
\[
  x^2+y^2+z^2=a^2
\]
的外侧，求
\[
  I=\iint_S x\,\mathrm{d}y\mathrm{d}z+y\,\mathrm{d}z\mathrm{d}x+z\,\mathrm{d}x\mathrm{d}y.
\]
\tcblower
令
\[
  \mathbf F=(x,y,z).
\]
则
\[
  \operatorname{div}\mathbf F=3.
\]
由 Gauss 公式，
\[
  I=\iiint_V\operatorname{div}\mathbf F\,\mathrm{d}V
  =3\iiint_V1\,\mathrm{d}V.
\]
球体体积为
\[
  \frac{4}{3}\pi a^3,
\]
所以
\[
  I=3\cdot\frac{4}{3}\pi a^3=4\pi a^3.
\]
\textbf{提示：}闭曲面的第二型曲面积分优先考虑 Gauss 公式。
\end{example}

\begin{example}{开曲面补面后用Gauss公式}{ex-gauss-cap}
设 $S$ 为上半球面
\[
  x^2+y^2+z^2=a^2,\qquad z\geq0,
\]
取上侧，计算
\[
  I=\iint_S x\,\mathrm{d}y\mathrm{d}z+y\,\mathrm{d}z\mathrm{d}x+z\,\mathrm{d}x\mathrm{d}y.
\]
\tcblower
令
\[
  \mathbf F=(x,y,z).
\]
用底面圆盘
\[
  S_0:z=0,\qquad x^2+y^2\leq a^2
\]
补成闭曲面。闭曲面的外侧在 $S_0$ 上为下侧。因为在 $S_0$ 上 $z=0$，且法向量为 $(0,0,-1)$，所以补面通量为 $0$。

由 Gauss 公式，
\[
  I=\iiint_V3\,\mathrm{d}V
  =3\cdot\frac12\cdot\frac{4}{3}\pi a^3
  =2\pi a^3.
\]
\textbf{提示：}开曲面不能直接用 Gauss 公式，必须先补成闭曲面。
\end{example}

\begin{example}{Stokes公式：空间闭曲线积分}{ex-stokes-triangle}
设 $L$ 为三角形
\[
  (1,0,0)\to(0,1,0)\to(0,0,1)\to(1,0,0),
\]
求
\[
  I=\int_L (y-z)\,\mathrm{d}x+(z-x)\,\mathrm{d}y+(x-y)\,\mathrm{d}z.
\]
\tcblower
令
\[
  \mathbf F=(y-z,z-x,x-y).
\]
则
\[
  \operatorname{curl}\mathbf F=(-2,-2,-2).
\]
三角形所在平面为
\[
  x+y+z=1.
\]
按题给方向，对应面积向量为
\[
  \frac12\bigl((0,1,0)-(1,0,0)\bigr)\times\bigl((0,0,1)-(1,0,0)\bigr)
  =
  \frac12(1,1,1).
\]
由 Stokes 公式，
\[
  I=\iint_S\operatorname{curl}\mathbf F\cdot \mathbf n\,\mathrm{d}S
  =
  (-2,-2,-2)\cdot\frac12(1,1,1)
  =
  -3.
\]
\textbf{提示：}Stokes 公式最容易错方向，曲线方向和曲面法向必须符合右手法则。
\end{example}



\section{高频易错点}
\begin{mistake}{有奇点仍直接套Green公式}{singularmistake}
Green公式要求 $P,Q$ 在整个闭区域上有连续偏导数。如果原点是奇点，而曲线围住原点，必须挖去小圆或选择不含奇点的内边界，否则结论可能完全错误。
\end{mistake}
\begin{mistake}{第二型曲面积分方向弄反}{orientationmistake}
曲面 $z=f(x,y)$ 取上侧时，法向量可取 $(-f_x,-f_y,1)$；取下侧时取相反方向。半球“上侧”与“外侧”有时一致，有时要结合曲面位置判断。
\end{mistake}
\examkey\ 2025回忆卷和中科大A2卷都集中考查Green公式、Gauss公式、补面、补洞、对称性与方向判断。

\chapter{含参变量积分}
\section{核心定义}
\begin{definition}{含参变量正常积分}{paramint}
设 $f(x,t)$ 在 $[a,b]\times T$ 上定义，
\[
F(t)=\int_a^b f(x,t)\dd x
\]
称为含参变量积分。核心问题是：$F(t)$ 的连续性、可导性、极限与积分号能否交换。
\end{definition}

\begin{definition}{含参变量反常积分的一致收敛}{improperparam}
设 $F(A,t)=\int_a^A f(x,t)\dd x$。若当 $A\to\infty$ 时，$F(A,t)$ 关于 $t\in T$ 一致收敛，则称 $\int_a^\infty f(x,t)\dd x$ 在 $T$ 上一致收敛。
\end{definition}

\section{关键定理}
\begin{theorem}{连续性与可导性定理}{paramtheorem}
若 $f(x,t)$ 在 $[a,b]\times T$ 上连续，则 $F(t)=\int_a^b f(x,t)\dd x$ 连续。若 $f_t(x,t)$ 存在且连续，则
\[
F'(t)=\int_a^b \frac{\partial f}{\partial t}(x,t)\dd x.
\]
若上下限也依赖于 $t$，即 $F(t)=\int_{\alpha(t)}^{\beta(t)}f(x,t)\dd x$，则
\[
F'(t)=f(\beta(t),t)\beta'(t)-f(\alpha(t),t)\alpha'(t)+\int_{\alpha(t)}^{\beta(t)}f_t(x,t)\dd x.
\]
\end{theorem}
连续性定理告诉我们满足这个条件就有\[
    \lim_{y\to y_0}\int_{a}^{b}f(x,y)\dd x=\int_{a}^{b}\lim_{y\to y_0}f(x,y)\dd x
\]既极限运算和积分运算可以交换。而可导性定理给了我们另一种计算积分的好办法。

\begin{theorem}{积分换序定理}{f}
    若\(f(x,y)\)在\([a,b] \times [c,d]\)上连续，则\[
        \int_{c}^{d}\dd y\int_{a}^{b}f(x,y)\dd x=\int_{a}^{b}\dd x\int_{c}^{d}f(x,y)\dd y
    \]
\end{theorem}






\begin{theorem}{含参反常积分的M判别法}{paramM}
若 $|f(x,t)|\le g(x)$ 对所有 $t\in T$ 成立，且 $\int_a^\infty g(x)\dd x$ 收敛，则 $\int_a^\infty f(x,t)\dd x$ 在 $T$ 上绝对一致收敛。
\end{theorem}

\section{典型解题方法}
\begin{methodbox}{积分号下求导模板}
\begin{enumerate}
\item 明确参数是哪一个，先写 $F(t)$。
\item 若是正常积分，检查连续性即可；若是反常积分，要额外检查一致收敛或找可积控制函数。
\item 对 $f_t$ 求导后，检查 $f_t$ 的连续性或一致可积控制。
\item 写出Leibniz公式，特别注意上下限求导带来的边界项。
\item 若题目要求极限，常将极限写成含参积分并用控制收敛或一致收敛交换极限。
\end{enumerate}
\end{methodbox}

\begin{example}{含参定积分的连续性}{ex-param-continuity}
设
\[
  F(x)=\int_0^1 \frac{\ln(1+xt)}{1+t^2}\,\mathrm{d}t,\qquad 0\leq x\leq 2.
\]
证明 $F(x)$ 在 $[0,2]$ 上连续。
\tcblower
令
\[
  f(x,t)=\frac{\ln(1+xt)}{1+t^2}.
\]
当 $(x,t)\in[0,2]\times[0,1]$ 时，$1+xt\geq1$，所以 $f(x,t)$ 在矩形区域上连续。由紧集上一致连续性，对任意 $\varepsilon>0$，存在 $\delta>0$，当 $|x_1-x_2|<\delta$ 时，对所有 $t\in[0,1]$ 有
\[
  |f(x_1,t)-f(x_2,t)|<\varepsilon.
\]
于是
\[
  |F(x_1)-F(x_2)|
  \leq \int_0^1 |f(x_1,t)-f(x_2,t)|\,\mathrm{d}t
  <\varepsilon.
\]
故 $F$ 在 $[0,2]$ 上连续。

\textbf{提示：}定区间、连续被积函数的含参积分通常直接给出连续函数。
\end{example}

\begin{example}{积分换序的重要性}{}
计算累次积分 $\displaystyle\int_{0}^{1} \mathrm{d}y \int_{y}^{1} \mathrm{e}^{-x^2} \, \mathrm{d}x$。
\tcblower
被积函数 $\mathrm{e}^{-x^2}$ 的原函数不是初等函数，若按原积分次序先对 $x$ 积分，无法使用牛顿-莱布尼茨公式直接计算，因此必须通过交换积分次序求解。

原积分对应的积分区域为 $Y$ 型区域：
$$
D: \quad 0 \le y \le 1,\quad y \le x \le 1
$$
该区域由直线 $y=x$、$x=1$ 与 $y=0$ 围成，将其改写为 $X$ 型区域：
$$
D: \quad 0 \le x \le 1,\quad 0 \le y \le x
$$

根据积分换序定理，将先对 $x$ 后对 $y$ 的累次积分，交换为先对 $y$ 后对 $x$ 的累次积分：
$$
\int_{0}^{1} \mathrm{d}y \int_{y}^{1} \mathrm{e}^{-x^2} \, \mathrm{d}x = \int_{0}^{1} \mathrm{d}x \int_{0}^{x} \mathrm{e}^{-x^2} \, \mathrm{d}y
$$

先计算内层对 $y$ 的积分（将 $x$ 视为常数）：
$$
\int_{0}^{x} \mathrm{e}^{-x^2} \, \mathrm{d}y = \mathrm{e}^{-x^2} \cdot \int_{0}^{x} \mathrm{d}y = x\mathrm{e}^{-x^2}
$$

再计算外层对 $x$ 的积分，凑微分求解：
$$
\begin{aligned}
\int_{0}^{1} x\mathrm{e}^{-x^2} \, \mathrm{d}x &= -\frac{1}{2}\int_{0}^{1} \mathrm{e}^{-x^2} \, \mathrm{d}(-x^2) \\
&= -\frac{1}{2}\left. \mathrm{e}^{-x^2} \right|_{0}^{1} \\
&= -\frac{1}{2}\left( \mathrm{e}^{-1} - 1 \right) \\
&= \frac{1}{2}\left( 1 - \frac{1}{\mathrm{e}} \right)
\end{aligned}
$$
\end{example}




\begin{example}{用参数求经典积分}{ex-param-log}
设 $a,b>0$，证明
\[
  \int_0^{+\infty}\frac{e^{-ax}-e^{-bx}}{x}\,\mathrm{d}x
  =
  \ln\frac ba.
\]
\tcblower
固定 $a>0$，令
\[
  F(b)=\int_0^{+\infty}\frac{e^{-ax}-e^{-bx}}{x}\,\mathrm{d}x.
\]
形式求导得
\[
  F'(b)=\int_0^{+\infty}e^{-bx}\,\mathrm{d}x=\frac1b.
\]
在 $b$ 取任意闭区间 $[\beta_1,\beta_2]\subset(0,+\infty)$ 时，求导后的被积函数 $e^{-bx}$ 可被 $e^{-\beta_1x}$ 控制，因此可在积分号下求导。于是
\[
  F(b)=\int_a^b \frac1s\,\mathrm{d}s=\ln b-\ln a=\ln\frac ba,
\]
其中用到
\[
  F(a)=0.
\]

\textbf{提示：}含参积分计算中常把较难积分看作参数函数，先求导，再积分还原。
\end{example}

\begin{example}{经典的问题}{al}
计算泊松积分 $\displaystyle I(\alpha) = \int_{0}^{\pi}\ln(1-2\alpha \cos x+\alpha^2)\,\dd x$。
\tcblower
被积函数含参变量 $\alpha$，我们通过积分号下求导（积分与求导换序）求解，这是积分换序定理的核心应用场景之一。该积分的原函数不是初等函数，无法直接用牛顿-莱布尼茨公式计算。

\textbf{1. 当 $|\alpha| < 1$ 时}
此时 $1-2\alpha \cos x+\alpha^2 > 0$ 恒成立，被积函数及其对 $\alpha$ 的偏导数在积分区域上连续，满足积分与求导换序的条件。

对 $I(\alpha)$ 关于 $\alpha$ 求导，交换求导与积分的次序：
$$
I'(\alpha) = \frac{\mathrm{d}}{\mathrm{d}\alpha}\int_{0}^{\pi}\ln(1-2\alpha \cos x+\alpha^2)\,\dd x
= \int_{0}^{\pi} \frac{\partial}{\partial\alpha}\ln(1-2\alpha \cos x+\alpha^2)\,\dd x
$$

计算偏导数：
$$
\frac{\partial}{\partial\alpha}\ln(1-2\alpha \cos x+\alpha^2) = \frac{2\alpha - 2\cos x}{1-2\alpha \cos x+\alpha^2}
$$

代入后对被积函数变形化简：
$$
I'(\alpha) = \int_{0}^{\pi} \frac{2\alpha - 2\cos x}{1-2\alpha \cos x+\alpha^2}\,\dd x
= \frac{1}{\alpha}\int_{0}^{\pi}\dd x + \frac{\alpha^2 - 1}{\alpha}\int_{0}^{\pi} \frac{1}{1-2\alpha \cos x+\alpha^2}\,\dd x
$$

利用万能代换可证得常用积分公式：
$$\int_{0}^{\pi} \frac{1}{1-2\alpha \cos x+\alpha^2}\,\dd x = \frac{\pi}{1-\alpha^2} \quad (|\alpha|<1)$$

将其代入得：
$$
I'(\alpha) = \frac{\pi}{\alpha} + \frac{\alpha^2 - 1}{\alpha}\cdot\frac{\pi}{1-\alpha^2} = \frac{\pi}{\alpha} - \frac{\pi}{\alpha} = 0
$$

$I'(\alpha)\equiv0$ 说明 $I(\alpha)$ 在 $|\alpha|<1$ 时为常数，代入特殊值 $\alpha=0$ 确定常数：
$$
I(0) = \int_{0}^{\pi}\ln 1\,\dd x = 0
$$

因此当 $|\alpha| \le 1$ 时，$\displaystyle I(\alpha) = 0$。

\textbf{2. 当 $|\alpha| > 1$ 时}
令 $\displaystyle\beta = \frac{1}{\alpha}$，则 $|\beta| < 1$，对被积函数做代数变形：
$$
\ln(1-2\alpha \cos x+\alpha^2) = \ln\left[\alpha^2\left(1 - 2\cdot\frac{1}{\alpha}\cos x + \frac{1}{\alpha^2}\right)\right]
= 2\ln|\alpha| + \ln(1-2\beta \cos x+\beta^2)
$$

两边同时积分得：
$$
I(\alpha) = \int_{0}^{\pi}2\ln|\alpha|\,\dd x + \int_{0}^{\pi}\ln(1-2\beta \cos x+\beta^2)\,\dd x
= 2\pi\ln|\alpha| + I(\beta)
$$

由 $|\beta|<1$ 时 $I(\beta)=0$，因此当 $|\alpha| > 1$ 时：
$$
I(\alpha) = 2\pi\ln|\alpha|
$$

综上，泊松积分的最终结果为：
$$
\int_{0}^{\pi}\ln(1-2\alpha \cos x+\alpha^2)\,\dd x =
\begin{cases}
0, & |\alpha| \le 1 \\
2\pi\ln|\alpha|, & |\alpha| > 1
\end{cases}
$$
\end{example}


\section{高频易错点}
\begin{mistake}{反常积分中随意交换极限}{parammistake}
正常积分在紧区间上有连续性时比较安全；反常积分涉及无限区间或无界被积函数，必须补充一致收敛或可积控制函数，否则积分号下求导和极限交换可能失败。
\end{mistake}

\chapter{Fourier级数}
\section{核心定义}
\begin{definition}{三角级数与Fourier系数}{fouriercoef}
设 $f$ 在 $[-\pi,\pi]$ 上可积，其Fourier系数定义为
\[
a_n=\frac1\pi\int_{-\pi}^{\pi}f(x)\cos nx\dd x\quad(n\ge0),\qquad
b_n=\frac1\pi\int_{-\pi}^{\pi}f(x)\sin nx\dd x\quad(n\ge1).
\]
形式级数
\[
\frac{a_0}{2}+\sum_{n=1}^{\infty}(a_n\cos nx+b_n\sin nx)
\]
称为 $f$ 的Fourier级数。
\end{definition}

\section{关键定理}
\begin{theorem}{三角函数正交性}{orth}
在 $[-\pi,\pi]$ 上，
\[
\int_{-\pi}^{\pi}\cos nx\cos mx\dd x=\begin{cases}0,&m\ne n,\\\pi,&m=n\ne0,\end{cases}
\]
\[
\int_{-\pi}^{\pi}\sin nx\sin mx\dd x=\begin{cases}0,&m\ne n,\\\pi,&m=n,\end{cases}
\quad
\int_{-\pi}^{\pi}\sin nx\cos mx\dd x=0.
\]
正交性是系数公式的来源。
\end{theorem}

\begin{theorem}{Dirichlet收敛定理}{dirifourier}
若 $f$ 在 $[-\pi,\pi]$ 上分段光滑，则其Fourier级数在每点 $x$ 收敛到
\[
\frac{f(x^{+})+f(x{-})}{2}.
\]
特别地，在连续点收敛到 $f(x)$；在跳跃间断点收敛到左右极限平均值。
\end{theorem}分段光滑的函数\(f(x)\)是它可以被fourier展开的充分条件。


\begin{theorem}{Parseval等式}{parseval}
若 $f$ 满足相应条件，则
\[
\frac1\pi\int_{-\pi}^{\pi}|f(x)|^2\dd x=\frac{a_0^2}{2}+\sum_{n=1}^{\infty}(a_n^2+b_n^2).
\]
该式常用于求 $\sum 1/n^2$、$\sum 1/(2n-1)^2$ 等数值。
\end{theorem}
\textbf{证明思路：}可以直接展开硬算，这里略去过程。
\section{典型解题方法}

\begin{methodbox}{Fourier展开四步法}
\begin{enumerate}
\item 判断区间：$[-\pi,\pi]$ 用完整Fourier级数；$[0,\pi]$ 常用余弦或正弦展开。
\item 判断奇偶：偶函数只有余弦项，奇函数只有正弦项。
\item 计算系数：分部积分是多项式乘三角函数的主要工具。
\item 写收敛值：连续点取函数值，间断点取左右极限均值，端点按周期延拓处理。
\end{enumerate}
\end{methodbox}

\begin{example}{半区间余弦展开}{ex-fourier-half-cos}
将
\[
  f(x)=x,\qquad 0\leq x\leq \pi
\]
展开为余弦级数。
\tcblower
余弦级数对应偶延拓：
\[
  f(x)\sim \frac{a_0}{2}+\sum_{n=1}^{\infty}a_n\cos nx,
\]
其中
\[
  a_0=\frac{2}{\pi}\int_0^\pi x\,\mathrm{d}x=\pi.
\]
对 $n\geq1$，
\[
  a_n=\frac{2}{\pi}\int_0^\pi x\cos nx\,\mathrm{d}x.
\]
分部积分得
\[
  \int_0^\pi x\cos nx\,\mathrm{d}x
  =
  \left.\frac{x\sin nx}{n}\right|_0^\pi
  -
  \frac{1}{n}\int_0^\pi\sin nx\,\mathrm{d}x.
\]
第一项为 $0$，且
\[
  \int_0^\pi\sin nx\,\mathrm{d}x
  =
  \frac{1-(-1)^n}{n}.
\]
所以
\[
  a_n=\frac{2}{\pi}\cdot \frac{(-1)^n-1}{n^2}.
\]
因此
\[
  x=\frac{\pi}{2}
  +
  \frac{2}{\pi}\sum_{n=1}^{\infty}
  \frac{(-1)^n-1}{n^2}\cos nx,
  \qquad 0\leq x\leq\pi.
\]

\textbf{提示：}半区间余弦展开的本质是把函数偶延拓到 $[-\pi,\pi]$。
\end{example}


\begin{example}{Fourier级数在端点处的收敛值}{ex-fourier-endpoint-value}
设
\[
f(x)=
\begin{cases}
0, & -\pi<x\leq 0,\\
x, & 0<x\leq \pi.
\end{cases}
\]
求函数 $f$ 的 Fourier 级数在 $x=\pi$ 处收敛到的值。
\tcblower
\textbf{解法一：用收敛定理直接判断。}

把 $f$ 看成以 $2\pi$ 为周期的函数。由于 $f$ 在 $(-\pi,\pi]$ 上分段光滑，所以由 Fourier 级数的 Dirichlet 收敛定理可知：在任意一点 $x_0$，Fourier 级数收敛到
\[
\frac{f(x_0^{+})+f(x_0^{-})}{2}.
\]

现在取 $x_0=\pi$。由定义可知
\[
f(\pi-0)=\lim_{x\to \pi^-}f(x)=\pi.
\]
但右极限要按照 $2\pi$ 周期延拓来理解。因为
\[
\pi+0
\]
右侧对应到基本区间左端
\[
-\pi+0,
\]
而当 $-\pi<x\leq 0$ 时，
\[
f(x)=0.
\]
所以
\[
f(\pi+0)=f(-\pi+0)=0.
\]
因此 Fourier 级数在 $x=\pi$ 处收敛到
\[
\frac{f(\pi-0)+f(\pi+0)}{2}
=
\frac{\pi+0}{2}
=
\frac{\pi}{2}.
\]

\[
\boxed{\text{Fourier级数在 }x=\pi\text{ 处收敛到 }\frac{\pi}{2}.}
\]

\textbf{解法二：硬算 Fourier 系数验证。}

Fourier 级数写成
\[
\frac{a_0}{2}+\sum_{n=1}^{\infty}\left(a_n\cos nx+b_n\sin nx\right),
\]
其中
\[
a_0=\frac{1}{\pi}\int_{-\pi}^{\pi}f(x)\,\mathrm{d}x,\qquad
a_n=\frac{1}{\pi}\int_{-\pi}^{\pi}f(x)\cos nx\,\mathrm{d}x,
\]
\[
b_n=\frac{1}{\pi}\int_{-\pi}^{\pi}f(x)\sin nx\,\mathrm{d}x.
\]

由于 $f(x)=0$ 在 $(-\pi,0]$ 上，$f(x)=x$ 在 $(0,\pi]$ 上，所以
\[
a_0=\frac{1}{\pi}\int_0^\pi x\,\mathrm{d}x
=
\frac{1}{\pi}\cdot \frac{\pi^2}{2}
=
\frac{\pi}{2}.
\]
因此
\[
\frac{a_0}{2}=\frac{\pi}{4}.
\]

下面计算 $a_n$：
\[
a_n=\frac{1}{\pi}\int_0^\pi x\cos nx\,\mathrm{d}x.
\]
分部积分，取
\[
u=x,\qquad \mathrm{d}v=\cos nx\,\mathrm{d}x,
\]
则
\[
\mathrm{d}u=\mathrm{d}x,\qquad v=\frac{\sin nx}{n}.
\]
所以
\[
\int_0^\pi x\cos nx\,\mathrm{d}x
=
\left.\frac{x\sin nx}{n}\right|_0^\pi
-
\frac{1}{n}\int_0^\pi \sin nx\,\mathrm{d}x.
\]
由于
\[
\sin n\pi=0,\qquad \sin 0=0,
\]
第一项为 $0$。又
\[
\int_0^\pi \sin nx\,\mathrm{d}x
=
\left.-\frac{\cos nx}{n}\right|_0^\pi
=
\frac{1-(-1)^n}{n}.
\]
因此
\[
\int_0^\pi x\cos nx\,\mathrm{d}x
=
-\frac{1}{n}\cdot \frac{1-(-1)^n}{n}
=
\frac{(-1)^n-1}{n^2}.
\]
于是
\[
a_n=\frac{(-1)^n-1}{\pi n^2}.
\]

再计算 $b_n$：
\[
b_n=\frac{1}{\pi}\int_0^\pi x\sin nx\,\mathrm{d}x.
\]
分部积分，取
\[
u=x,\qquad \mathrm{d}v=\sin nx\,\mathrm{d}x,
\]
则
\[
\mathrm{d}u=\mathrm{d}x,\qquad v=-\frac{\cos nx}{n}.
\]
所以
\[
\int_0^\pi x\sin nx\,\mathrm{d}x
=
\left.-\frac{x\cos nx}{n}\right|_0^\pi
+
\frac{1}{n}\int_0^\pi \cos nx\,\mathrm{d}x.
\]
其中
\[
\left.-\frac{x\cos nx}{n}\right|_0^\pi
=
-\frac{\pi\cos n\pi}{n}
=
-\frac{\pi(-1)^n}{n},
\]
且
\[
\frac{1}{n}\int_0^\pi \cos nx\,\mathrm{d}x
=
\left.\frac{\sin nx}{n^2}\right|_0^\pi
=
0.
\]
故
\[
\int_0^\pi x\sin nx\,\mathrm{d}x
=
-\frac{\pi(-1)^n}{n}.
\]
因此
\[
b_n=\frac{(-1)^{n+1}}{n}.
\]

所以 Fourier 级数为
\[
\frac{\pi}{4}
+
\sum_{n=1}^{\infty}
\left[
\frac{(-1)^n-1}{\pi n^2}\cos nx
+
\frac{(-1)^{n+1}}{n}\sin nx
\right].
\]

令 $x=\pi$，则
\[
\cos n\pi=(-1)^n,\qquad \sin n\pi=0.
\]
因此级数在 $x=\pi$ 处的形式值为
\[
\frac{\pi}{4}
+
\sum_{n=1}^{\infty}
\frac{(-1)^n-1}{\pi n^2}(-1)^n.
\]
化简求和项：
\[
\frac{(-1)^n-1}{\pi n^2}(-1)^n
=
\frac{1-(-1)^n}{\pi n^2}.
\]
所以
\[
\frac{\pi}{4}
+
\sum_{n=1}^{\infty}
\frac{1-(-1)^n}{\pi n^2}.
\]
当 $n$ 为偶数时，
\[
1-(-1)^n=0;
\]
当 $n$ 为奇数时，
\[
1-(-1)^n=2.
\]
故
\[
\sum_{n=1}^{\infty}
\frac{1-(-1)^n}{\pi n^2}
=
\frac{2}{\pi}\sum_{k=0}^{\infty}\frac{1}{(2k+1)^2}.
\]
又
\[
\sum_{k=0}^{\infty}\frac{1}{(2k+1)^2}
=
\sum_{n=1}^{\infty}\frac1{n^2}
-
\sum_{k=1}^{\infty}\frac1{(2k)^2}
=
\frac{\pi^2}{6}-\frac14\cdot\frac{\pi^2}{6}
=
\frac{\pi^2}{8}.
\]
因此
\[
\frac{2}{\pi}\sum_{k=0}^{\infty}\frac{1}{(2k+1)^2}
=
\frac{2}{\pi}\cdot\frac{\pi^2}{8}
=
\frac{\pi}{4}.
\]
最终得到
\[
\frac{\pi}{4}+\frac{\pi}{4}
=
\frac{\pi}{2}.
\]

\[
\boxed{\text{答案为 }\frac{\pi}{2}.}
\]

\textbf{易错点：}
题目问的是 Fourier 级数在 $x=\pi$ 处的收敛值，不是函数值 $f(\pi)$。虽然
\[
f(\pi)=\pi,
\]
但由于周期延拓后 $x=\pi$ 是跳跃点，Fourier 级数收敛到左右极限的平均值：
\[
\frac{\pi+0}{2}=\frac{\pi}{2}.
\]
\end{example}

\begin{example}{利用 Parseval 等式求级数和}{ex-fourier-parseval-odd}
利用 Parseval 等式求
\[
  \sum_{k=0}^{\infty}\frac{1}{(2k+1)^2}.
\]
\tcblower
取函数
\[
  f(x)=
  \begin{cases}
  -1, & -\pi<x<0,\\
  1, & 0<x<\pi.
  \end{cases}
\]
它的 Fourier 系数为
\[
  a_0=0,\qquad a_n=0,\qquad
  b_{2k}=0,\qquad b_{2k+1}=\frac{4}{\pi(2k+1)}.
\]
由 Parseval 等式，
\[
  \frac{1}{\pi}\int_{-\pi}^{\pi}|f(x)|^2\,\mathrm{d}x
  =
  \sum_{n=1}^{\infty}b_n^2.
\]
左边为
\[
  \frac{1}{\pi}\int_{-\pi}^{\pi}1\,\mathrm{d}x=2.
\]
右边为
\[
  \sum_{k=0}^{\infty}
  \left(\frac{4}{\pi(2k+1)}\right)^2
  =
  \frac{16}{\pi^2}
  \sum_{k=0}^{\infty}\frac{1}{(2k+1)^2}.
\]
因此
\[
  2=
  \frac{16}{\pi^2}
  \sum_{k=0}^{\infty}\frac{1}{(2k+1)^2}.
\]
所以
\[
  \sum_{k=0}^{\infty}\frac{1}{(2k+1)^2}
  =
  \frac{\pi^2}{8}.
\]

\textbf{提示：}Parseval 等式常用来求平方倒数型级数，尤其是只含奇数项或交错项的级数。这个简单的级数事实上也可以通过巴塞尔问题既著名的\(\sum_{k=1}^{\infty}\frac{1}{k^2}=\frac{\pi^2}{6}\)转化得到，这里不再赘述。
\end{example}



\section{高频易错点}
\begin{mistake}{忘记端点按周期延拓}{fouriermistake}
半区间展开在端点处的和不一定等于原函数端点值，而是等于延拓后左右极限平均。做求和题时，代入 $x=0$、$x=\pi$ 前必须判断延拓方式。
\end{mistake}



\part{附录：三套试卷原题与极详细解析}
\appendix

\chapter{浙江大学数学分析（甲）II（H）2024--2025春夏期末回忆卷}
\section*{原题重排}
\begin{problem}{A卷原题}
\begin{enumerate}[label=\textbf{A\arabic*.}]
\item 叙述二元函数 $f(x,y)$ 在 $(x_0,y_0)$ 可微的定义，并证明
\[
f(x,y)=\begin{cases}\dfrac{x^2y}{\sqrt{x^2+y^2}},&(x,y)\ne(0,0),\\0,&(x,y)=(0,0)
\end{cases}
\]
在 $(0,0)$ 处可微。
\item 计算：
\begin{enumerate}[label=(\arabic*)]
\item 求由 $x^2+y^2+z\ee^z=2$，$x^2+xy+y^2=1$ 确定的空间曲线在 $(1,-1,0)$ 处的切线；
\item 求 $I_1=\iint_D |y-x^2|\dd x\dd y$，其中 $D=[-1,1]\times[0,2]$；
\item 对曲线 $L:y=\sqrt{1-x^2}$，方向从 $(1,0)$ 到 $(-1,0)$，求
\[
I_2=\int_L(-2x\ee^{-x^2}\sin y-y)\dd x+\ee^{-x^2}\cos y\dd y;
\]
\item 对上半球面 $S:x^2+y^2+z^2=1,z\ge0$，取上侧，求
\[
I_3=\iint_S(x^2-x)\dd y\dd z+(y^2-y)\dd z\dd x+(z^2+1)\dd x\dd y;
\]
\item 求
\[
I_4=\int_{-1}^{1}\dd x\int_0^{\sqrt{1-x^2}}\dd y\int_1^{1+\sqrt{1-x^2-y^2}}\frac{\dd z}{\sqrt{x^2+y^2+z^2}}.
\]
\end{enumerate}
\item 设 $f(x,y)$ 定义在 $D=(0,1)\times(0,1)$ 上，$f_y'(x,y)$ 有界，且固定 $y$ 时 $f(x,y)$ 关于 $x$ 连续。证明 $f$ 在 $D$ 上连续。
\item 对 $f(x,y,z)=2x^2+2xy+2y^2-3z^2$，方向 $\vec l=(1,-1,0)$，点 $P$ 在 $S:x^2+2y^2+3z^2=1$ 上，求 $\dfrac{\partial f}{\partial \vec l}(P)$ 的最大值。
\item 对函数 $f(x)=1+x$，$x\in[0,\pi]$，有 $f(x)=\dfrac{a_0}{2}+\sum_{n=1}^{\infty}a_n\cos nx$，求 $\lim_{n\to\infty}n^2\sin a_{2n-1}$。
\item 函数列 $\{f_n\}$ 满足 $f_1(x)=f(x)$ 为 $[0,1]$ 上连续函数，且 $f_{n+1}(x)=\int_x^1 f_n(t)\dd t$。证明函数项级数 $\sum_{n=1}^{\infty}f_n(x)$ 在 $[0,1]$ 上一致收敛。
\item 设 $a_n>0$ 且 $\sum a_n$ 收敛，记 $R_n=\sum_{k=n+1}^{\infty}a_k$。证明：
\begin{enumerate}[label=(\arabic*)]
\item $\sum_{n=1}^{\infty}\dfrac{a_n}{R_{n-1}}$ 发散；
\item 对任意 $p>0$，$\sum_{n=1}^{\infty}\dfrac{a_n}{R_{n-1}^{1-p}}$ 收敛。
\end{enumerate}
\end{enumerate}
\end{problem}

\section*{逐题极详细解析}
\begin{solutionbox}{A1 可微定义与原点可微证明}

\thinking\ 证明某个分段函数在原点可微，最稳妥的方法不是先依赖偏导连续，因为偏导往往不连续；而是直接按定义验证。先算原点偏导，确定线性主部，再证明剩余项是 $o(r)$，其中 $r=\sqrt{x^2+y^2}$。

\textbf{第一步：写出可微定义。} 二元函数 $f$ 在 $(x_0,y_0)$ 可微，是指存在常数 $A,B$，使得当 $h\to0,k\to0$ 时
\[
f(x_0+h,y_0+k)-f(x_0,y_0)=Ah+Bk+o\left(\sqrt{h^2+k^2}\right).
\]
若可微，则 $A=f_x(x_0,y_0)$，$B=f_y(x_0,y_0)$。

\textbf{第二步：求原点偏导。} 由定义
\[
f_x(0,0)=\lim_{h\to0}\frac{f(h,0)-f(0,0)}{h}=\lim_{h\to0}\frac{0}{h}=0,
\]
因为 $y=0$ 时分子 $x^2y=0$。同理
\[
f_y(0,0)=\lim_{k\to0}\frac{f(0,k)-f(0,0)}{k}=0.
\]
所以若可微，其线性主部只能是 $0\cdot x+0\cdot y=0$。

\textbf{第三步：按定义估计余项。} 记 $r=\sqrt{x^2+y^2}$。当 $(x,y)\ne(0,0)$ 时，
\[
|f(x,y)|=\frac{|x|^2|y|}{\sqrt{x^2+y^2}}=\frac{|x|^2|y|}{r}.
\]
由于 $|x|\le r$，$|y|\le r$，所以
\[
|x|^2|y|\le r^3,
\qquad
|f(x,y)|\le r^2.
\]
于是
\[
\frac{|f(x,y)-f(0,0)-0x-0y|}{\sqrt{x^2+y^2}}
=\frac{|f(x,y)|}{r}\le r\to0.
\]
这说明余项为 $o(r)$，故 $f$ 在 $(0,0)$ 处可微，且全微分为 $\dd f(0,0)=0$。

\checklist\ 本题的关键不是“偏导存在”，而是“函数值比距离高一阶小”。沿任何路径都被 $r^2$ 控制，因此可微成立。
\end{solutionbox}

\begin{solutionbox}{A2(1) 隐式空间曲线切线}
\point\ 空间曲线由两个曲面交线给出；切向量等于两个法向量的叉乘。

设
\[
F_1(x,y,z)=x^2+y^2+z\ee^z-2,
\qquad
F_2(x,y,z)=x^2+xy+y^2-1.
\]
曲线是 $F_1=0$ 与 $F_2=0$ 的交线。曲面 $F_i=0$ 在一点的法向量为 $\grad F_i$，因此交线切向量垂直于两个法向量，可取
\[
\bm v=\grad F_1(P)\times \grad F_2(P).
\]
计算梯度：
\[
\grad F_1=(2x,2y,\ee^z+z\ee^z),
\qquad
\grad F_2=(2x+y,x+2y,0).
\]
在 $P=(1,-1,0)$ 处，
\[
\grad F_1(P)=(2,-2,1),
\qquad
\grad F_2(P)=(1,-1,0).
\]
叉乘为
\[
\begin{vmatrix}
\bm i&\bm j&\bm k\\
2&-2&1\\
1&-1&0
\end{vmatrix}
=(1,1,0).
\]
故切线为
\[
\boxed{(x,y,z)=(1,-1,0)+t(1,1,0),\quad t\in\R.}
\]
\thinking\ 交线切线题不要尝试直接解出参数方程。两个约束曲面的梯度给出两个法向量，叉乘就是切向量。
\end{solutionbox}

\begin{solutionbox}{A2(2) 含绝对值二重积分}
\point\ 矩形区域上绝对值积分，按符号分界线 $y=x^2$ 分段。

固定 $x\in[-1,1]$，令 $a=x^2$，则 $0\le a\le1$。在 $0\le y\le a$ 上，$|y-a|=a-y$；在 $a\le y\le2$ 上，$|y-a|=y-a$。因此
\[
\int_0^2|y-x^2|\dd y
=\int_0^{x^2}(x^2-y)\dd y+\int_{x^2}^{2}(y-x^2)\dd y.
\]
逐项计算：
\[
\int_0^{x^2}(x^2-y)\dd y=x^2\cdot x^2-\frac{x^4}{2}=\frac{x^4}{2},
\]
\[
\int_{x^2}^{2}(y-x^2)\dd y=\frac{4-x^4}{2}-x^2(2-x^2)=2-2x^2+\frac{x^4}{2}.
\]
相加得
\[
\int_0^2|y-x^2|\dd y=x^4-2x^2+2.
\]
故
\[
I_1=\int_{-1}^{1}(x^4-2x^2+2)\dd x
=2\int_0^1(x^4-2x^2+2)\dd x
=2\left(\frac15-\frac23+2\right)=\boxed{\frac{46}{15}}.
\]
\checklist\ 含绝对值的积分第一步一定是找零点，本题零点是 $y=x^2$，且它完全落在积分区间 $[0,2]$ 内。
\end{solutionbox}

\begin{solutionbox}{A2(3) 第二类曲线积分与精确微分}
\point\ 计算第二类曲线积分，这里考虑用Green公式。

设曲线\(l_1\)为题目所述曲线，\(l_2\)为\((-1,0)\to(1,0)\)的线段。那么由Green公式就有\[
    \int_{l_1}\cdots+\int_{l_2}\cdots=\iint_{D}(\frac{\partial Q}{\partial x}-\frac{\partial P}{\partial y})\dd x\dd y
\]这里区域D就是\(l_1,l_2\)围成的区域。计算\[
\frac{\partial Q}{\partial x}-\frac{\partial P}{\partial y}=1
\]故\begin{align*}
    I+\int_{l_2}\cdots&=\iint_{D}\dd x\dd y\\
 \implies I&=\frac{\pi}{2}
\end{align*}

\thinking\ 计算第二型曲线积分常有定义法，Green公式，Stokes公式等。
\end{solutionbox}

\begin{solutionbox}{A2(4) 上半球通量}
\point\ 第二型曲面积分、Gauss公式、补面方向。

令
\[
\bm F=(x^2-x,\ y^2-y,\ z^2+1).
\]
所求即上半球 $S$ 上取上侧的通量 $\iint_S\bm F\cdot\bm n\dd S$。上半球上侧即外侧。用底面圆盘 $D_0:x^2+y^2\le1,z=0$ 补成闭曲面 $\Sigma=S\cup D_0$，整体取外侧。由Gauss公式
\[
\iint_\Sigma \bm F\cdot\bm n\dd S=\iiint_V\diver \bm F\dd V,
\]
其中 $V$ 是上半球体。散度为
\[
\diver\bm F=2x+2y+2z.
\]
由于上半球体关于 $x$、$y$ 对称，$\iiint_V2x\dd V=\iiint_V2y\dd V=0$。故闭曲面通量为
\[
\iiint_V2z\dd V.
\]
柱坐标下 $0\le r\le1$，$0\le\theta\le2\pi$，$0\le z\le\sqrt{1-r^2}$，于是
\[
\iiint_V2z\dd V=\int_0^{2\pi}\int_0^1\int_0^{\sqrt{1-r^2}}2z\,r\dd z\dd r\dd\theta
=2\pi\int_0^1 r(1-r^2)\dd r=\frac{\pi}{2}.
\]
底面 $D_0$ 的外法向为 $(0,0,-1)$，在 $z=0$ 上 $\bm F\cdot(0,0,-1)=-(z^2+1)=-1$，故底面通量
\[
\iint_{D_0}-1\dd A=-\pi.
\]
因此
\[
I_3+(-\pi)=\frac{\pi}{2},
\qquad
\boxed{I_3=\frac{3\pi}{2}.}
\]
\checklist\ 补面法必须核对补面方向。这里底面取闭曲面外侧，是向下，不是向上。
\end{solutionbox}

\begin{solutionbox}{A2(5) 三重积分的几何换元}
\point\ 空间区域识别、球坐标、积分限。

原积分区域满足
\[
-1\le x\le1,
\quad 0\le y\le\sqrt{1-x^2},
\quad 1\le z\le1+\sqrt{1-x^2-y^2}.
\]
前两个条件说明 $(x,y)$ 在单位上半圆盘内；第三个条件说明 $z$ 位于平面 $z=1$ 与球
\[
x^2+y^2+(z-1)^2=1
\]
的上半部分之间。因此区域是以 $(0,0,1)$ 为球心、半径 $1$ 的球的上半部分，同时 $y\ge0$。

被积函数为 $1/\sqrt{x^2+y^2+z^2}$，适合以原点为球心的球坐标：
\[
x=\rho\sin\varphi\cos\theta,
\quad y=\rho\sin\varphi\sin\theta,
\quad z=\rho\cos\varphi.
\]
条件 $y\ge0$ 给出 $0\le\theta\le\pi$。条件 $z\ge1$ 给出 $\rho\cos\varphi\ge1$。球面
\[
x^2+y^2+(z-1)^2=1
\]
化为
\[
\rho^2-2\rho\cos\varphi+1=1,
\quad\text{即}\quad \rho=2\cos\varphi.
\]
同时 $z\ge1$ 要求 $\rho\ge\sec\varphi$。为了有非空区间，需要
\[
\sec\varphi\le2\cos\varphi
\Longleftrightarrow
\cos^2\varphi\ge\frac12
\Longleftrightarrow
0\le\varphi\le\frac\pi4.
\]
故
\[
I_4=\int_0^\pi\int_0^{\pi/4}\int_{\sec\varphi}^{2\cos\varphi}\frac1\rho \rho^2\sin\varphi\dd\rho\dd\varphi\dd\theta.
\]
先对 $\rho$ 积分：
\[
\int_{\sec\varphi}^{2\cos\varphi}\rho\dd\rho
=\frac12(4\cos^2\varphi-\sec^2\varphi).
\]
于是
\[
I_4=\frac\pi2\int_0^{\pi/4}(4\cos^2\varphi-\sec^2\varphi)\sin\varphi\dd\varphi.
\]
分别计算
\[
\int_0^{\pi/4}4\cos^2\varphi\sin\varphi\dd\varphi
=\frac43\left(1-\frac{1}{2\sqrt2}\right),
\]
\[
\int_0^{\pi/4}\sec^2\varphi\sin\varphi\dd\varphi
=\int_0^{\pi/4}\tan\varphi\sec\varphi\dd\varphi
=\sqrt2-1.
\]
代入化简得
\[
I_4=\boxed{\frac{\pi(7-4\sqrt2)}{6}}.
\]
\end{solutionbox}

\begin{solutionbox}{A3 连续性证明}
\point\ 偏导有界推出对一个变量Lipschitz，另一个变量用已知连续性。

设 $|f_y'(x,y)|\le M$。任取 $(x_0,y_0)\in D$，证明 $f(x,y)\to f(x_0,y_0)$。将差值拆成两段：
\[
|f(x,y)-f(x_0,y_0)|\le |f(x,y)-f(x,y_0)|+|f(x,y_0)-f(x_0,y_0)|.
\]
对第一项，固定 $x$，对一元函数 $y\mapsto f(x,y)$ 用中值定理：存在介于 $y$ 与 $y_0$ 之间的 $\eta$，使
\[
f(x,y)-f(x,y_0)=f_y'(x,\eta)(y-y_0),
\]
所以
\[
|f(x,y)-f(x,y_0)|\le M|y-y_0|.
\]
对第二项，由题设固定 $y_0$ 时 $x\mapsto f(x,y_0)$ 连续，故当 $x\to x_0$ 时
\[
|f(x,y_0)-f(x_0,y_0)|\to0.
\]
因此给定 $\varepsilon>0$，先令 $|y-y_0|<\varepsilon/(2M)$ 控制第一项，再令 $|x-x_0|$ 足够小控制第二项，即得 $f$ 在 $(x_0,y_0)$ 连续。由于点任意，$f$ 在 $D$ 上连续。

\thinking\ 本题的本质是“一个方向有一致的Lipschitz控制，另一个方向逐点连续”，两段路径连接即可。
\end{solutionbox}

\begin{solutionbox}{A4 椭球上方向导数最大值}
\point\ 方向导数、线性函数在椭球上的最大值。

方向 $\vec l=(1,-1,0)$ 的单位向量为
\[
\bm u=\frac{1}{\sqrt2}(1,-1,0).
\]
函数梯度为
\[
\grad f=(4x+2y,\ 2x+4y,\ -6z).
\]
方向导数为
\[
D_{\bm u}f=\grad f\cdot\bm u
=\frac{(4x+2y)-(2x+4y)}{\sqrt2}
=\sqrt2(x-y).
\]
因此问题化为在椭球
\[
x^2+2y^2+3z^2=1
\]
上求 $x-y$ 的最大值。

用Cauchy不等式写
\[
x-y=1\cdot x+\left(-\frac1{\sqrt2}\right)(\sqrt2y)+0\cdot(\sqrt3z).
\]
于是
\[
(x-y)^2\le \left(1^2+\frac12+0\right)(x^2+2y^2+3z^2)=\frac32.
\]
故 $x-y\le\sqrt{3/2}$，从而
\[
D_{\bm u}f\le \sqrt2\sqrt{\frac32}=\sqrt3.
\]
等号可由Cauchy等号条件达到，因此最大值为
\[
\boxed{\sqrt3}.
\]
\checklist\ 若方向导数定义要求方向向量单位化，必须先除以 $\sqrt2$。这一步漏掉会差一个倍数。
\end{solutionbox}

\begin{solutionbox}{A5 Fourier余弦系数极限}
\point\ 半区间余弦系数、分部积分、等价无穷小。

在 $[0,\pi]$ 上的余弦展开系数为
\[
a_n=\frac2\pi\int_0^\pi(1+x)\cos nx\dd x.
\]
由于
\[
\int_0^\pi \cos nx\dd x=\frac{\sin n\pi}{n}=0,
\]
只需计算 $\int_0^\pi x\cos nx\dd x$。分部积分：取 $u=x$，$\dd v=\cos nx\dd x$，则 $\dd u=\dd x$，$v=\sin nx/n$，
\[
\int_0^\pi x\cos nx\dd x=\left.\frac{x\sin nx}{n}\right|_0^\pi-\frac1n\int_0^\pi\sin nx\dd x.
\]
第一项为 $0$，第二项为
\[
-\frac1n\left[-\frac{\cos nx}{n}\right]_0^\pi
=\frac{\cos n\pi-1}{n^2}=\frac{(-1)^n-1}{n^2}.
\]
因此
\[
a_n=\frac2\pi\frac{(-1)^n-1}{n^2}.
\]
当指标为 $2n-1$ 时，$(-1)^{2n-1}=-1$，故
\[
a_{2n-1}= -\frac{4}{\pi(2n-1)^2}.
\]
由于 $a_{2n-1}\to0$，$\sin a_{2n-1}\sim a_{2n-1}$，所以
\[
\lim_{n\to\infty}n^2\sin a_{2n-1}
=\lim_{n\to\infty}n^2\left(-\frac{4}{\pi(2n-1)^2}\right)
=\boxed{-\frac1\pi}.
\]
\end{solutionbox}

\begin{solutionbox}{A6 递推积分函数项级数的一致收敛}
这个问题有一定难度，我们考虑用最大值直接放缩，然后看看会不会有一些好结果。

设
\[
M=\max_{0\le x\le1}|f(x)|.
\]
由于 $f$ 连续，$M<\infty$。我们证明归纳估计
\[
|f_n(x)|\le M\frac{(1-x)^{n-1}}{(n-1)!},\qquad x\in[0,1].
\]
当 $n=1$ 时，右端为 $M$，成立。若对 $n$ 成立，则
\[
|f_{n+1}(x)|=\left|\int_x^1 f_n(t)\dd t\right|
\le\int_x^1 |f_n(t)|\dd t
\le \int_x^1 M\frac{(1-t)^{n-1}}{(n-1)!}\dd t.
\]
计算积分：
\[
\int_x^1(1-t)^{n-1}\dd t=\frac{(1-x)^n}{n}.
\]
故
\[
|f_{n+1}(x)|\le M\frac{(1-x)^n}{n!}.
\]
归纳完成。

又因 $0\le1-x\le1$，
\[
|f_n(x)|\le \frac{M}{(n-1)!}.
\]
数项级数 $\sum_{n=1}^{\infty}\frac{M}{(n-1)!}=Me$ 收敛。由Weierstrass判别法，函数项级数 $\sum f_n(x)$ 在 $[0,1]$ 上绝对一致收敛。

\thinking\ 递推积分题常会产生阶乘。每积分一次，幂次升一、阶乘多一，这是统一估计的来源。
\end{solutionbox}

\begin{solutionbox}{A7 正项级数余项和构造题}

因为 $a_n>0$ 且 $\sum a_n$ 收敛，所以尾和
\[
R_n=\sum_{k=n+1}^{\infty}a_k
\]
满足 $R_n>0$ 且 $R_n\downarrow0$。并且
\[
R_{n-1}=a_n+R_n,
\qquad
\frac{R_n}{R_{n-1}}=1-\frac{a_n}{R_{n-1}}.
\]

\textbf{(1) 证明发散。} 记 $b_n=a_n/R_{n-1}$，则 $0<b_n<1$，且
\[
\prod_{k=1}^{n}(1-b_k)=\prod_{k=1}^{n}\frac{R_k}{R_{k-1}}=\frac{R_n}{R_0}\to0.
\]
若 $\sum b_n$ 收敛，则由于 $0<b_n<1$，无穷乘积 $\prod(1-b_n)$ 应收敛到正数；这与上式趋于 $0$ 矛盾。因此
\[
\sum_{n=1}^{\infty}\frac{a_n}{R_{n-1}}
\]
发散。

\textbf{(2) 证明收敛。} 要证
\[
\sum a_nR_{n-1}^{p-1}
\]
收敛。

若 $p\ge1$，则 $R_{n-1}^{p-1}\le R_0^{p-1}$，从而
\[
a_nR_{n-1}^{p-1}\le R_0^{p-1}a_n,
\]
右端级数收敛，故结论成立。

若 $0<p<1$，函数 $t^{p-1}$ 在 $(0,\infty)$ 上递减。由于 $R_n<R_{n-1}$，对 $t\in[R_n,R_{n-1}]$ 有 $t^{p-1}\ge R_{n-1}^{p-1}$，因此
\[
\int_{R_n}^{R_{n-1}}t^{p-1}\dd t\ge (R_{n-1}-R_n)R_{n-1}^{p-1}=a_nR_{n-1}^{p-1}.
\]
于是
\[
\sum_{n=1}^{N}a_nR_{n-1}^{p-1}
\le\sum_{n=1}^{N}\int_{R_n}^{R_{n-1}}t^{p-1}\dd t
=\int_{R_N}^{R_0}t^{p-1}\dd t
\le \frac{R_0^p}{p}.
\]
部分和有界且各项非负，所以级数收敛。

\thinking\ 第一问考虑转化为无穷乘积；第二问用积分把 $R_{n-1}-R_n=a_n$ 变成望远镜。这是正项级数构造题的典型技巧。
\end{solutionbox}

\chapter{中国科学技术大学数学分析A2期末资料}
\section*{原题重排}
\begin{problem}{B卷原题}
\begin{enumerate}[label=\textbf{B\arabic*.}]
\item 设 $\alpha>1$。计算
\[
\iint_D\frac{y^\alpha}{\sqrt{x+y^2}}\dd x\dd y,
\]
其中 $D$ 由 $x=0$，$y=0$，$x+y^2=1$ 围成。
\item 计算：
\begin{enumerate}[label=(\arabic*)]
\item 设 $L:|x|^3+3|y|=1$ 为封闭曲线，计算 $\int_L|x|^3\dd s$；
\item 计算 $\oint_L\dfrac{-y}{4x^2+y^2}\dd x+\dfrac{x}{4x^2+y^2}\dd y$，其中 $L:x^2+y^2=9$ 取逆时针方向。
\end{enumerate}
\item 求
\[
I=\iint_S(x^2-x)\dd y\dd z+(y^2-y)\dd z\dd x+(z^2-z)\dd x\dd y,
\]
其中 $S$ 是半径为 $R$ 的上半球面 $x^2+y^2+z^2=R^2,z\ge0$，方向朝上。
\item 设 $\Sigma:x^2/a^2+y^2/b^2+z^2/c^2=1$，法向朝外。计算
\[
I=\iint_\Sigma x^n\dd y\dd z+y^n\dd z\dd x+z^n\dd x\dd y.
\]
\item 设平面向量场
\[
\bm F=\left(-\frac{y}{a^2x^2+b^2y^2},\frac{x}{a^2x^2+b^2y^2}\right),\quad a,b>0.
\]
(1) 求 $\bm F$ 在 $D_1=\{(x,y):x>0\}$ 上的所有势函数；(2) 证明 $\bm F$ 不是 $D_2=\R^2\setminus\{(0,0)\}$ 上的保守场。
\item 设 $D=\{(x,y):x^2+y^2\le R^2\}$，$f\in C^2(D)$ 且 $\Delta f=x+x^2+y^2$。计算 $\iint_Df(x,y)\dd x\dd y$。
\item 设 $f\in C[0,1]$，$D=\{0\le\theta\le\pi/2,0\le\varphi\le\pi/2\}$。证明
\[
\iint_D(\sin\theta+\cos\varphi)f(1-\sin\theta\cos\varphi)\dd\theta\dd\varphi
=\pi\int_0^1f(x)\dd x.
\]
\end{enumerate}
\end{problem}

\section*{逐题极详细解析}
\begin{solutionbox}{B1 特殊换元二重积分}

\thinking\ 被积函数中有 $\sqrt{x+y^2}$，边界也有 $x+y^2=1$，所以我们希望令 $x+y^2=r^2$。同时 $x\ge0,y\ge0$，自然可取
\[
x=r^2\cos^2\varphi,
\qquad
y=r\sin\varphi,
\qquad
0\le\varphi\le\frac\pi2.
\]
此时
\[
x+y^2=r^2\cos^2\varphi+r^2\sin^2\varphi=r^2.
\]
边界 $x+y^2=1$ 化为 $r=1$，坐标轴 $x=0,y=0$ 对应 $\varphi=\pi/2$ 与 $\varphi=0$，故区域变为矩形
\[
0\le r\le1,
\qquad
0\le\varphi\le\frac\pi2.
\]
计算Jacobian：
\[
\frac{\partial(x,y)}{\partial(r,\varphi)}=
\begin{vmatrix}
2r\cos^2\varphi&-2r^2\cos\varphi\sin\varphi\\
\sin\varphi&r\cos\varphi
\end{vmatrix}
=2r^2\cos\varphi.
\]
因此
\[
\begin{aligned}
I&=\int_0^1\int_0^{\pi/2}
\frac{(r\sin\varphi)^\alpha}{r}\cdot 2r^2\cos\varphi\dd\varphi\dd r\\
&=2\left(\int_0^1r^{\alpha+1}\dd r\right)
\left(\int_0^{\pi/2}\sin^\alpha\varphi\cos\varphi\dd\varphi\right)\\
&=2\cdot\frac1{\alpha+2}\cdot\frac1{\alpha+1}.
\end{aligned}
\]
故
\[
\boxed{I=\frac{2}{(\alpha+1)(\alpha+2)}}.
\]
\checklist\ 本题最容易错的是把边界看错。换元的目的正是让 $x+y^2$ 变成 $r^2$，区域也随之变成标准矩形。
\end{solutionbox}

\begin{solutionbox}{B2(1) 第一型曲线积分}
\point\ 对称性、弧长微元。

曲线 $|x|^3+3|y|=1$ 关于两坐标轴对称，被积函数 $|x|^3$ 和弧长元也关于两轴对称，所以积分等于第一象限部分的 $4$ 倍。第一象限中
\[
y=\frac{1-x^3}{3},\quad 0\le x\le1.
\]
于是
\[
y'(x)=-x^2,
\qquad
\dd s=\sqrt{1+x^4}\dd x.
\]
故
\[
\int_L|x|^3\dd s=4\int_0^1x^3\sqrt{1+x^4}\dd x.
\]
令 $t=x^4$，则 $\dd t=4x^3\dd x$，得到
\[
4\int_0^1x^3\sqrt{1+x^4}\dd x=
\int_0^1\sqrt{1+t}\dd t
=\left.\frac{2}{3}(1+t)^{3/2}\right|_0^1.
\]
因此
\[
\boxed{\int_L|x|^3\dd s=\frac23(2\sqrt2-1)}.
\]
\end{solutionbox}

\begin{solutionbox}{B2(2) 绕奇点场的第二型曲线积分}
\point\ Green公式补洞、奇点处理、方向。

向量场
\[
P=-\frac{y}{4x^2+y^2},
\qquad
Q=\frac{x}{4x^2+y^2}
\]
在原点无定义，而圆 $x^2+y^2=9$ 围住原点，所以不能直接在圆盘上套Green公式。

取内边界椭圆 $L_1:4x^2+y^2=4$，使其包围原点。圆 $L$ 与椭圆 $L_1$ 之间的环域没有奇点，且可直接验证 $Q_x-P_y=0$。环域的正向边界为外边界 $L$ 逆时针加内边界 $L_1$ 顺时针。因此
\[
\oint_LP\dd x+Q\dd y+\oint_{-L_1}P\dd x+Q\dd y=0,
\]
从而
\[
\oint_LP\dd x+Q\dd y=\oint_{L_1}P\dd x+Q\dd y
\]
其中右侧 $L_1$ 取逆时针方向。

参数化 $L_1$：
\[
x=\cos t,
\qquad
y=2\sin t,
\qquad
0\le t\le2\pi.
\]
则
\[
\dd x=-\sin t\dd t,
\qquad
\dd y=2\cos t\dd t,
\qquad
4x^2+y^2=4.
\]
代入：
\[
P\dd x+Q\dd y=\frac{-2\sin t}{4}(-\sin t\dd t)+\frac{\cos t}{4}(2\cos t\dd t)
=\frac12\dd t.
\]
故
\[
\boxed{\oint_LP\dd x+Q\dd y=\int_0^{2\pi}\frac12\dd t=\pi.}
\]
\thinking\ 对含 $1/(x^2+y^2)$ 型奇点的场，核心不是算偏导，而是先判断区域是否含奇点。含奇点则补洞。
\end{solutionbox}

\begin{solutionbox}{B3 上半球第二型曲面积分}
\point\ 曲面积分通量形式、球面法向、对称性。

令
\[
\bm F=(x^2-x,\ y^2-y,\ z^2-z).
\]
上半球面方向朝上，与外法向一致。半径为 $R$ 的球面上外单位法向为
\[
\bm n=\frac{1}{R}(x,y,z).
\]
因此
\[
I=\frac1R\iint_S[(x^2-x)x+(y^2-y)y+(z^2-z)z]\dd S.
\]
即
\[
I=\frac1R\iint_S(x^3+y^3+z^3-x^2-y^2-z^2)\dd S.
\]
球面上 $x^2+y^2+z^2=R^2$，故
\[
I=\frac1R\iint_S(x^3+y^3+z^3-R^2)\dd S.
\]
由关于 $yz$ 面和 $xz$ 面的对称性，
\[
\iint_Sx^3\dd S=0,
\qquad
\iint_Sy^3\dd S=0.
\]
于是
\[
I=\frac1R\iint_S z^3\dd S-R\iint_S\dd S.
\]
上半球面积为 $2\pi R^2$，第二项为 $2\pi R^3$。再用球坐标
\[
z=R\cos\theta,
\quad
\dd S=R^2\sin\theta\dd\theta\dd\varphi,
\quad
0\le\theta\le\frac\pi2,
\quad
0\le\varphi\le2\pi.
\]
则
\[
\frac1R\iint_Sz^3\dd S
=\frac1R\int_0^{2\pi}\int_0^{\pi/2}R^3\cos^3\theta\,R^2\sin\theta\dd\theta\dd\varphi
=2\pi R^4\cdot\frac14=\frac{\pi R^4}{2}.
\]
故
\[
\boxed{I=\frac{\pi R^4}{2}-2\pi R^3.}
\]
\end{solutionbox}

\begin{solutionbox}{B4 椭球通量积分}
\point\ Gauss公式、线性伸缩换元、奇偶性。

将积分看成向量场 $\bm F=(x^n,y^n,z^n)$ 穿过椭球面的外通量。由Gauss公式，
\[
I=\iiint_V\diver\bm F\dd V
=n\iiint_V(x^{n-1}+y^{n-1}+z^{n-1})\dd V,
\]
其中 $V$ 是椭球内部。

作线性换元
\[
x=au,
\qquad y=bv,
\qquad z=cw.
\]
椭球 $V$ 变为单位球 $B:u^2+v^2+w^2\le1$，Jacobian为 $abc$。因此
\[
I=nabc\iiint_B(a^{n-1}u^{n-1}+b^{n-1}v^{n-1}+c^{n-1}w^{n-1})\dd u\dd v\dd w.
\]
若 $n$ 为偶数，则 $n-1$ 为奇数，$u^{n-1},v^{n-1},w^{n-1}$ 在对称单位球上积分均为 $0$，故 $I=0$。

若 $n$ 为奇数，则 $n-1$ 为偶数，三个积分相等：
\[
\iiint_Bu^{n-1}\dd V=\iiint_Bv^{n-1}\dd V=\iiint_Bw^{n-1}\dd V.
\]
用球坐标算
\[
\iiint_Bw^{n-1}\dd V
=\int_0^1r^{n+1}\dd r\int_0^\pi\cos^{n-1}\theta\sin\theta\dd\theta\int_0^{2\pi}\dd\varphi.
\]
由于 $n-1$ 为偶数，
\[
\int_0^\pi\cos^{n-1}\theta\sin\theta\dd\theta=\frac2n,
\qquad
\int_0^1r^{n+1}\dd r=\frac1{n+2}.
\]
故
\[
\iiint_Bw^{n-1}\dd V=\frac{4\pi}{n(n+2)}.
\]
代回得
\[
\boxed{
I=\begin{cases}
0,&n\text{为偶数},\\[2mm]
\dfrac{4\pi}{n+2}abc(a^{n-1}+b^{n-1}+c^{n-1}),&n\text{为奇数}.
\end{cases}}
\]
\checklist\ 若原题限定 $n$ 为正奇数，则只需写第二行。若题面只说正整数，必须讨论奇偶。
\end{solutionbox}

\begin{solutionbox}{B5 势函数与非保守场}
\point\ 势函数构造、路径无关、单连通性。

\textbf{(1) 求右半平面上的势函数。} 右半平面 $x>0$ 单连通且不含原点。取基点 $(1,0)$ 到 $(u,v)$ 的折线路径：先沿 $x$ 轴到 $(u,0)$，再竖直到 $(u,v)$。第一段 $y=0$，$\dd y=0$，且 $P=0$，积分为 $0$。第二段 $x=u$，$\dd x=0$，所以
\[
\phi(u,v)=\int_0^v\frac{u}{a^2u^2+b^2y^2}\dd y.
\]
计算积分：
\[
\int\frac{u}{a^2u^2+b^2y^2}\dd y
=\frac1{ab}\arctan\frac{by}{au}.
\]
因此所有势函数为
\[
\boxed{\phi(x,y)=\frac1{ab}\arctan\frac{by}{ax}+C,
\quad x>0.}
\]
可直接求偏导核对 $\phi_x=P,\phi_y=Q$。

\textbf{(2) 证明穿孔平面上不是保守场。} 若在 $D_2=\R^2\setminus\{0\}$ 上是保守场，则任意闭曲线积分应为 $0$。取椭圆
\[
L:a^2x^2+b^2y^2=1
\]
逆时针方向。参数化
\[
x=\frac1a\cos t,
\quad y=\frac1b\sin t,
\quad 0\le t\le2\pi.
\]
则 $a^2x^2+b^2y^2=1$，
\[
\dd x=-\frac1a\sin t\dd t,
\quad
\dd y=\frac1b\cos t\dd t.
\]
代入积分得
\[
P\dd x+Q\dd y
=\left(-\frac{\sin t/b}{1}\right)\left(-\frac{\sin t}{a}\dd t\right)
+\left(\frac{\cos t/a}{1}\right)\left(\frac{\cos t}{b}\dd t\right)
=\frac1{ab}\dd t.
\]
故
\[
\oint_L\bm F\cdot\dd\bm r=\frac1{ab}\int_0^{2\pi}\dd t=\frac{2\pi}{ab}\ne0.
\]
所以 $\bm F$ 不是 $D_2$ 上的保守场。

\thinking\ 局部满足 $P_y=Q_x$ 不代表全局保守；区域是否单连通是关键。
\end{solutionbox}

\begin{solutionbox}{B6 Poisson方程与调和函数平均值}

题设
\[
\Delta f=x+x^2+y^2.
\]
先构造一个特解 $g$，使 $\Delta g=x+x^2+y^2$。取
\[
g(x,y)=\frac{x^3}{6}+\frac{x^4}{12}+\frac{y^4}{12}.
\]
因为
\[
\frac{\partial^2}{\partial x^2}\frac{x^3}{6}=x,
\quad
\frac{\partial^2}{\partial x^2}\frac{x^4}{12}=x^2,
\quad
\frac{\partial^2}{\partial y^2}\frac{y^4}{12}=y^2,
\]
所以 $\Delta g=x+x^2+y^2$。令
\[
h=f-g.
\]
则 $\Delta h=0$，即 $h$ 在圆盘上调和。调和函数在圆盘上满足面积平均值性质：
\[
\iint_Dh\dd A=\pi R^2h(0,0).
\]
由于 $g(0,0)=0$，所以 $h(0,0)=f(0,0)$。

接着计算 $g$ 的积分。奇函数项 $x^3/6$ 在圆盘上积分为 $0$。又由对称性
\[
\iint_Dx^4\dd A=\iint_Dy^4\dd A.
\]
极坐标中
\[
\iint_Dx^4\dd A=\int_0^R\int_0^{2\pi}r^4\cos^4\theta\,r\dd\theta\dd r
=\frac{R^6}{6}\cdot\frac{3\pi}{4}=\frac{\pi R^6}{8}.
\]
因此
\[
\iint_Dg\dd A=\frac1{12}\left(\frac{\pi R^6}{8}+\frac{\pi R^6}{8}\right)=\frac{\pi R^6}{48}.
\]
最终
\[
\boxed{\iint_Df\dd A=\pi R^2f(0,0)+\frac{\pi R^6}{48}.}
\]
\checklist\ 此题不能从 $\Delta f$ 直接恢复 $f$；必须“特解 + 调和函数”分解。
\end{solutionbox}

\begin{solutionbox}{B7 二重积分恒等式证明}
我们考虑重积分换元,分成两个部分。\\
要证明
\[
I=\iint_D(\sin\theta+\cos\varphi)f(1-\sin\theta\cos\varphi)\dd\theta\dd\varphi
=\pi\int_0^1f(x)\dd x.
\]
把 $I$ 分成两部分 $I=I_1+I_2$。

\textbf{第一部分。} 对
\[
I_1=\iint_D\sin\theta\,f(1-\sin\theta\cos\varphi)\dd\theta\dd\varphi
\]
令
\[
u=\sin\theta\cos\varphi,
\quad
v=\cos\theta.
\]
在第一象限参数区域内，该换元把区域对应到单位四分圆 $u\ge0,v\ge0,u^2+v^2\le1$ 的适当区域；Jacobian会把 $\sin\theta\dd\theta\dd\varphi$ 化为平面面积元的一维投影积分。等价地，固定 $u=\sin\theta\cos\varphi$ 后，垂直截长给出
\[
I_1=\frac\pi2\int_0^1f(1-u)\dd u=\frac\pi2\int_0^1f(x)\dd x.
\]

\textbf{第二部分。} 类似地对
\[
I_2=\iint_D\cos\varphi\,f(1-\sin\theta\cos\varphi)\dd\theta\dd\varphi
\]
作对称换元，可得同样结果
\[
I_2=\frac\pi2\int_0^1f(x)\dd x.
\]
两式相加得到结论。

\end{solutionbox}

\chapter{南开数学分析II 2023--2024学年第二学期期末A卷}
\section*{原题重排}
\begin{problem}{C卷原题}
\begin{enumerate}[label=\textbf{C\arabic*.}]
\item 计算二重积分
\[
\iint_D\frac{xy+1}{x^2+y^2+3}\dd x\dd y,
\quad
D=\{(x,y)\in\R^2:x^2+y^2\le1,\ y\ge0\}.
\]
\item 设
\[
f(x,y)=\begin{cases}\dfrac{xy}{\sqrt{x^2+y^2}},&x^2+y^2\ne0,\\0,&x^2+y^2=0.
\end{cases}
\]
证明 $f$ 在原点连续，偏导数 $f_x'(x,y),f_y'(x,y)$ 在 $\R^2$ 上处处存在且有界，但 $f$ 在原点不可微。
\item 求曲面 $(x^2+y^2+z^2)^3=a^2z^4$ 所围区域体积，其中 $a>0$。
\item 设 $f(u,v)$ 在 $\R^2$ 上二次连续可微，且 $f_{uu}+f_{vv}=1$。令
\[
g(x,y)=f\left(xy,\frac{x^2-y^2}{2}\right).
\]
证明 $g_{xx}+g_{yy}=x^2+y^2$。
\item 求函数 $f(x,y)=xy$ 在闭区域
\[
D=\{(x,y)\in\R^2:x+y\le3,\\\ x^2+y^2\le8\}
\]
上的最大值和最小值。
\item 设 $0<a<b$，$f\in C[a,b]$，且 $\int_a^bf(x)\dd x=0$。证明：对任意正实数 $p$，存在 $\xi\in(a,b)$，使
\[
\int_a^\xi f(x)\dd x=p\xi f(\xi).
\]
\item 设 $f(x,t)$ 在 $[0,1]\times[0,1]$ 上二次连续可微，对任意 $(x,t)$ 有 $f_t=f_{xx}$ 且 $|f_x|\le1$。
\begin{enumerate}[label=(\arabic*)]
\item 证明对任意 $x,y,t_1,t_2\in[0,1]$，有
\[
\left|\int_x^y f(u,t_1)\dd u-\int_x^y f(u,t_2)\dd u\right|\le2|t_1-t_2|.
\]
\item 证明对任意 $x,t_1,t_2\in[0,1]$，有
\[
|f(x,t_1)-f(x,t_2)|\le5|t_1-t_2|^{1/2}.
\]
\end{enumerate}
\end{enumerate}
\end{problem}

\section*{逐题极详细解析}
\begin{solutionbox}{C1 半圆盘二重积分}

区域 $D$ 是上半单位圆盘，关于 $y$ 轴对称。被积函数拆为
\[
\frac{xy+1}{x^2+y^2+3}=\frac{xy}{x^2+y^2+3}+\frac1{x^2+y^2+3}.
\]
第一项关于 $x$ 是奇函数，而区域关于 $x\mapsto -x$ 对称，因此
\[
\iint_D\frac{xy}{x^2+y^2+3}\dd x\dd y=0.
\]
第二项用极坐标：$x=r\cos\theta,y=r\sin\theta$，其中 $0\le r\le1$，$0\le\theta\le\pi$。于是
\[
I=\int_0^\pi\int_0^1\frac{r}{r^2+3}\dd r\dd\theta.
\]
内层积分为
\[
\int_0^1\frac{r}{r^2+3}\dd r
=\frac12\ln(r^2+3)\bigg|_0^1
=\frac12\ln\frac43.
\]
所以
\[
\boxed{I=\frac\pi2\ln\frac43.}
\]
\thinking\ 二重积分先看对称性，能消掉奇函数部分会显著减少计算。
\end{solutionbox}

\begin{solutionbox}{C2 连续、偏导有界但不可微}
\point\ 多元函数可微定义、偏导存在与可微的区别。

\textbf{连续性。} 对 $(x,y)\ne(0,0)$，令 $r=\sqrt{x^2+y^2}$，则
\[
|f(x,y)|=\frac{|xy|}{r}\le\frac{x^2+y^2}{2r}=\frac r2\to0.
\]
故 $f(x,y)\to0=f(0,0)$，所以 $f$ 在原点连续。

\textbf{原点偏导数。}
\[
f_x(0,0)=\lim_{h\to0}\frac{f(h,0)-0}{h}=0,
\qquad
f_y(0,0)=\lim_{k\to0}\frac{f(0,k)-0}{k}=0.
\]

\textbf{非原点偏导。} 当 $(x,y)\ne(0,0)$ 时，
\[
f(x,y)=xy(x^2+y^2)^{-1/2}.
\]
对 $x$ 求导：
\[
f_x=y(x^2+y^2)^{-1/2}+xy\left(-\frac12\right)(x^2+y^2)^{-3/2}\cdot2x.
\]
化简得
\[
f_x=\frac{y(x^2+y^2)-x^2y}{(x^2+y^2)^{3/2}}=\frac{y^3}{(x^2+y^2)^{3/2}}.
\]
同理
\[
f_y=\frac{x^3}{(x^2+y^2)^{3/2}}.
\]
于是
\[
|f_x|\le1,
\qquad |f_y|\le1
\]
在非原点成立，原点也为 $0$，所以偏导处处存在且有界。

\textbf{不可微。} 若 $f$ 在原点可微，由原点偏导均为 $0$，线性主部为 $0$，应有
\[
\frac{f(x,y)}{\sqrt{x^2+y^2}}\to0.
\]
但沿 $y=x$，
\[
f(x,x)=\frac{x^2}{\sqrt{2x^2}}=\frac{|x|}{\sqrt2},
\quad
\sqrt{x^2+y^2}=\sqrt2|x|,
\]
故
\[
\frac{f(x,x)}{\sqrt{x^2+x^2}}=\frac12,
\]
不趋于 $0$。矛盾，因此不可微。

\checklist\ 偏导有界仍不能推出可微；偏导连续才是常用充分条件。
\end{solutionbox}

\begin{solutionbox}{C3 曲面围成区域体积}

考虑使用球坐标
\[
x=\rho\sin\varphi\cos\theta,
\quad y=\rho\sin\varphi\sin\theta,
\quad z=\rho\cos\varphi.
\]
于是
\[
x^2+y^2+z^2=\rho^2,
\qquad z=\rho\cos\varphi.
\]
边界方程
\[
(x^2+y^2+z^2)^3=a^2z^4
\]
化为
\[
\rho^6=a^2\rho^4\cos^4\varphi.
\]
除去原点后得
\[
\rho^2=a^2\cos^4\varphi,
\quad
\rho=a\cos^2\varphi.
\]
由于 $\cos^2\varphi\ge0$，全空间角度都覆盖，故
\[
0\le\theta\le2\pi,
\quad
0\le\varphi\le\pi,
\quad
0\le\rho\le a\cos^2\varphi.
\]
体积为
\[
V=\int_0^{2\pi}\int_0^\pi\int_0^{a\cos^2\varphi}\rho^2\sin\varphi\dd\rho\dd\varphi\dd\theta.
\]
先积分 $\rho$：
\[
\int_0^{a\cos^2\varphi}\rho^2\dd\rho=\frac{a^3\cos^6\varphi}{3}.
\]
于是
\[
V=\frac{2\pi a^3}{3}\int_0^\pi\cos^6\varphi\sin\varphi\dd\varphi.
\]
令 $u=\cos\varphi$，$\dd u=-\sin\varphi\dd\varphi$，上下限从 $1$ 到 $-1$，
\[
\int_0^\pi\cos^6\varphi\sin\varphi\dd\varphi=\int_{-1}^{1}u^6\dd u=\frac{2}{7}.
\]
因此
\[
\boxed{V=\frac{4\pi a^3}{21}.}
\]
\end{solutionbox}

\begin{solutionbox}{C4 Laplace算子在特殊变换下的链式计算}。

令
\[
u=xy,
\qquad
v=\frac{x^2-y^2}{2},
\qquad
g(x,y)=f(u,v).
\]
先计算一阶导数：
\[
u_x=y,
\quad u_y=x,
\quad v_x=x,
\quad v_y=-y.
\]
二阶导数为
\[
u_{xx}=u_{yy}=0,
\quad v_{xx}=1,
\quad v_{yy}=-1.
\]
对复合函数用二阶链式法则：
\[
\begin{aligned}
g_{xx}&=f_{uu}u_x^2+2f_{uv}u_xv_x+f_{vv}v_x^2+f_u u_{xx}+f_vv_{xx},\\
g_{yy}&=f_{uu}u_y^2+2f_{uv}u_yv_y+f_{vv}v_y^2+f_u u_{yy}+f_vv_{yy}.
\end{aligned}
\]
相加得
\[
\begin{aligned}
g_{xx}+g_{yy}
&=f_{uu}(u_x^2+u_y^2)+2f_{uv}(u_xv_x+u_yv_y)+f_{vv}(v_x^2+v_y^2)\\
&\quad +f_u(u_{xx}+u_{yy})+f_v(v_{xx}+v_{yy}).
\end{aligned}
\]
逐项计算：
\[
u_x^2+u_y^2=y^2+x^2,
\quad
v_x^2+v_y^2=x^2+y^2,
\]
\[
u_xv_x+u_yv_y=yx+x(-y)=0,
\]
\[
u_{xx}+u_{yy}=0,
\quad
v_{xx}+v_{yy}=1-1=0.
\]
因此
\[
g_{xx}+g_{yy}=(x^2+y^2)(f_{uu}+f_{vv}).
\]
由题设 $f_{uu}+f_{vv}=1$，得
\[
\boxed{g_{xx}+g_{yy}=x^2+y^2.}
\]
\thinking\ 本题若直接展开容易乱。按“二阶链式法则 - 交叉项 - Laplace项”分组，结构会非常清楚。
\end{solutionbox}

\begin{solutionbox}{C5 闭区域上 $xy$ 的最大最小值}

区域
\[
D=\{x+y\le3,\ x^2+y^2\le8\}
\]
是闭有界集，$xy$ 连续，因此最大值和最小值存在。

\textbf{内部临界点。} 在内部，
\[
\grad(xy)=(y,x)=0
\]
给出 $(0,0)$，函数值为 $0$。

\textbf{圆边界。} 在 $x^2+y^2=8$ 上，令 $s=x+y$，则
\[
xy=\frac{(x+y)^2-(x^2+y^2)}2=\frac{s^2-8}{2}.
\]
同时还要满足 $s\le3$。圆上 $s$ 的范围是 $[-4,4]$，加上 $s\le3$ 后为 $[-4,3]$。因此圆弧上
\[
xy=\frac{s^2-8}{2},\quad -4\le s\le3.
\]
最大值在 $|s|$ 最大处 $s=-4$ 取得，值
\[
\frac{16-8}{2}=4,
\]
对应点 $(-2,-2)$。最小值在 $s=0$ 取得，值
\[
\frac{0-8}{2}=-4,
\]
对应点 $(2,-2)$ 与 $(-2,2)$，它们满足 $x+y=0\le3$。

\textbf{直线边界。} 在 $x+y=3$ 且 $x^2+y^2\le8$ 上，令 $y=3-x$，则
\[
xy=x(3-x)=-x^2+3x.
\]
直线与圆交点由
\[
x^2+(3-x)^2=8
\]
给出
\[
2x^2-6x+1=0,
\quad
x=\frac{3\pm\sqrt7}{2}.
\]
该线段上二次函数最大值在 $x=3/2$ 处，为 $9/4$，端点值较小。与圆边界得到的 $4$ 和 $-4$ 比较后，全局最大值为 $4$，最小值为 $-4$。

故
\[
\boxed{\max_Dxy=4,
\quad
\min_Dxy=-4.}
\]
\checklist\ 闭区域极值题必须检查所有边界，不可只用Lagrange乘数得到部分点。
\end{solutionbox}

\begin{solutionbox}{C6 Rolle定理构造辅助函数}

令
\[
F(x)=\int_a^x f(t)\dd t.
\]
则 $F$ 在 $[a,b]$ 上连续、在 $(a,b)$ 上可导，且
\[
F'(x)=f(x),
\quad
F(a)=0,
\quad
F(b)=\int_a^bf(t)\dd t=0.
\]
题目要证明存在 $\xi$，使
\[
F(\xi)=p\xi f(\xi)=p\xi F'(\xi).
\]
等价于
\[
F'(\xi)-\frac{1}{p\xi}F(\xi)=0.
\]
这提示构造一个函数，使其导数正好含有上式。因为 $a>0$，可定义
\[
H(x)=\frac{F(x)}{x^{1/p}}.
\]
则 $H$ 在 $[a,b]$ 连续，在 $(a,b)$ 可导，且
\[
H(a)=\frac{F(a)}{a^{1/p}}=0,
\quad
H(b)=\frac{F(b)}{b^{1/p}}=0.
\]
由Rolle定理，存在 $\xi\in(a,b)$，使 $H'(\xi)=0$。

计算导数：
\[
H'(x)=x^{-1/p}F'(x)-\frac1p x^{-1/p-1}F(x).
\]
令 $x=\xi$，并用 $H'(\xi)=0$，得
\[
\xi^{-1/p}f(\xi)-\frac1p\xi^{-1/p-1}F(\xi)=0.
\]
两边乘以 $p\xi^{1/p+1}$，得
\[
F(\xi)=p\xi f(\xi).
\]
即
\[
\boxed{\int_a^\xi f(x)\dd x=p\xi f(\xi).}
\]
\thinking\ 看到 $F=p x F'$ 这类式子，应想到通过 $F/x^{1/p}$ 的导数构造Rolle定理。
\end{solutionbox}

\begin{solutionbox}{C7}
略
\end{solutionbox}

\backmatter
\chapter*{后记}
\addcontentsline{toc}{chapter}{后记}
\qquad 回头翻看这几十页的内容，从最开始零散的草稿笔记，到一点点补全例题、标注易错点、调整格式，不知不觉竟攒成了一本完整的讲义。整理的过程对我自己而言，也是一次重新梳理、重新理解的过程，收获远比我最开始预想的多。\par
我知道学数学的日子，很多时候是孤独的,或者说自己一个人努力，尝试变好的过程是孤独的，这些时刻我都经历过。我也曾经觉得，「把数学学好」「去参加竞赛」「去让更多人知道你」，这些都是别人的事，是我这种起点的人不敢想的事。但走着走着我就发现，很多你现在觉得「我肯定做不到」的事，其实都没那么遥远。\par
追溯最初的起点，如果高中没有想提前自学数学，如果高考后没有想着报数学专业，如果在曾经的专业没有坚定自己的信念想转来这里，如果我挂科了之后就彻底自暴自弃，那么也就不会有你看到这段文字，也不会有我能被更多人认识，但是那么多的如果出现在我面前我都选择坚持下来，我认为这和你的天资无关，这只和你的内心信念有关，既然我这种天资愚笨的人尚能坚持自己做出这些，你又怎么可能不行呢。\par
我觉得我自己只是一个敢想敢做的普通人，我还有很多事情想做，希望在未来都能做到。\par
最后我想说，写下这些的时候我只是一个普通的大二学生，没什么了不起的。如果这本我熬了很多夜写出来的讲义，能帮你少走一点点弯路，能在你觉得学不动的时候给你一点点底气，甚至能让你某一刻突然觉得「好像我也能学好数学」—— 那对我而言，就是莫大的荣幸。感谢。\par
祝大家期末顺利，都能拿到自己满意的成绩，也祝你能活出自己想要的样子。
\vfill

\begin{flushright}
Symphony\\
2026 年 6 月于武汉科技大学
\end{flushright}


\end{document}
